\chapter{Examples}
\label{chap:pet-examples}

This chapter computes curvature for warped products, Lie-group metrics, and
Riemannian submersions. The statements below use the curvature conventions fixed
in Chapter 3.

\section{Computational Simplifications}

The first three propositions reduce sectional and Ricci curvature computations to
diagonal curvature data in an orthonormal frame.

\begin{proposition}[sectional curvature bounds from a diagonalized curvature operator]
  \label{prop:pet-ch4-curvature-operator-diagonal-bounds}
  \lean{PetersenLib.secBoundsFromDiagonalCurvatureOperator}
  \source{petersen:page-0133}
  \leanok
  \dcref{ch4:4.1.1}
  Let $e_i$ be an orthonormal basis for $T_pM$. If $e_i\wedge e_j$ diagonalize the
  curvature operator, $\mathfrak{R}(e_i\wedge e_j)=\lambda_{ij}e_i\wedge e_j$, then for any
  plane $\pi\subset T_pM$, $\mathrm{sec}(\pi)\in[\min\lambda_{ij},\max\lambda_{ij}]$.
\end{proposition}
\begin{proof}
  \uses{prop:pet-ch4-curvature-operator-diagonal-bounds}
  \leanok
  \sketch
  For an orthonormal basis $v,w$ of $\pi$, expand the unit bivector
  $v\wedge w=\sum_{i<j}a_{ij}e_i\wedge e_j$. Self-adjointness and
  $\sum_{i<j}a_{ij}^2=1$ give
  $\operatorname{sec}(\pi)=\sum_{i<j}\lambda_{ij}a_{ij}^2$, which lies between
  the least and greatest eigenvalues.
\end{proof}

\begin{proposition}[pairwise vanishing diagonalizes the curvature operator]
  \label{prop:pet-ch4-triple-vanishing-diagonalizes-operator}
  \lean{PetersenLib.diagonalCurvatureOperatorFromTripleVanishing}
  \source{petersen:page-0134}
  \leanok
  \dcref{ch4:4.1.2}
  Let $e_i$ be an orthonormal basis for $T_pM$. If $\mathfrak{R}(e_i,e_j)e_k=0$ whenever
  $i,j,k$ are mutually distinct, then $e_i\wedge e_j$ diagonalize the curvature
  operator.
\end{proposition}
\begin{proof}
  \uses{prop:pet-ch4-triple-vanishing-diagonalizes-operator}
  \leanok
  \sketch
  Curvature symmetries give
  \[
    g(\mathfrak R(e_i\wedge e_j),e_k\wedge e_\ell)
    =-g(\mathfrak R(e_i,e_j)e_k,e_\ell)
    =g(\mathfrak R(e_i,e_j)e_\ell,e_k).
  \]
  The hypothesis makes this zero unless $\{k,\ell\}=\{i,j\}$; hence every
  off-diagonal matrix entry of $\mathfrak R$ in the indicated bivector basis
  vanishes.
\end{proof}

\begin{proposition}[a criterion for diagonalizing the Ricci tensor]
  \label{prop:pet-ch4-ricci-diagonal-criterion}
  \lean{PetersenLib.ricciDiagonalFromTripleVanishing}
  \source{petersen:page-0134}
  \leanok
  \dcref{ch4:4.1.3}
  Let $e_i$ be an orthonormal basis for $T_pM$. If $g(\mathfrak{R}(e_i,e_j)e_k,e_i)=0$
  whenever three of the indices $i,j,k$ are mutually distinct, then $e_i$
  diagonalize $\mathrm{Ric}$.
\end{proposition}
\begin{proof}
  \uses{prop:pet-ch4-ricci-diagonal-criterion}
  \leanok
  \sketch
  For $i\ne j$,
  $g(\operatorname{Ric}(e_i),e_j)=\sum_k g(\mathfrak R(e_i,e_k)e_k,e_j)$.
  Terms with $k\notin\{i,j\}$ vanish by hypothesis, while those with $k=i$ or
  $k=j$ vanish by the alternating curvature symmetries. Thus every off-diagonal
  Ricci component is zero.
\end{proof}

\section{Warped Products}

The first calculations derive the sphere, rotationally symmetric, and doubly warped
product curvature operators from the Euclidean model.

\subsection{Spheres}

The Euclidean distance Hessian determines the induced curvature of round spheres.

\begin{proposition}[Hessian of the Euclidean distance function]
  \label{prop:pet-ch4-hess-r-euclidean}
  \lean{PetersenLib.hessDistanceEuclidean}
  \source{petersen:page-0135}
  On $\mathbb{R}^n$, write the polar-coordinate metric $g=dr^2+g_r$, $g_r=r^2ds_{n-1}^2$,
  where $ds_{n-1}^2$ is the canonical metric on $S^{n-1}(1)$ and $r(x)=|x|$. Then
  $\mathrm{Hess}\,r=\tfrac1r g_r$.
\end{proposition}
\begin{proof}
  \uses{prop:pet-ch4-hess-r-euclidean}
  \sketch
  Both $dr$ and $ds_{n-1}^2$ are invariant under the radial flow. Therefore
  \[
    2\operatorname{Hess}r=L_{\partial_r}g
    =L_{\partial_r}(r^2ds_{n-1}^2)
    =2r\,ds_{n-1}^2=\frac2r g_r.
  \]
\end{proof}

\begin{proposition}[constant curvature of Euclidean spheres]
  \label{prop:pet-ch4-euclidean-sphere-constant-curvature}
  \lean{PetersenLib.sphereConstantCurvatureFromEuclideanTangential}
  \source{petersen:page-0135}
  \uses{prop:pet-ch4-hess-r-euclidean, prop:pet-ch4-triple-vanishing-diagonalizes-operator}
  For $n\ge3$, the level sphere $S^{n-1}(R)\subset\mathbb{R}^n$ (with its induced metric
  $g_R=R^2ds_{n-1}^2$) has constant curvature $R^{-2}$.
\end{proposition}
\begin{proof}
  \uses{prop:pet-ch4-euclidean-sphere-constant-curvature}
  \sketch
  By \cref{prop:pet-ch4-hess-r-euclidean}, the second fundamental form of the
  level sphere is $\mathrm{II}=r^{-1}g_r$. Since the ambient curvature is zero,
  the Gauss equation gives
  \[
    \mathfrak R^\top(X,Y)Z
    =r^{-2}\bigl(g_r(Y,Z)X-g_r(X,Z)Y\bigr).
  \]
  Thus \cref{prop:pet-ch4-triple-vanishing-diagonalizes-operator} applies and
  every bivector eigenvalue, hence every sectional curvature, equals $r^{-2}$.
  Put $r=R$.
\end{proof}

\subsection{Product Spheres}

Product connections split by factor, so pure and mixed bivectors have separate
curvature eigenvalues.

\begin{definition}[product sphere metric $S^n_a\times S^m_b$]
  \label{def:pet-ch4-product-sphere-metric}
  \lean{PetersenLib.productSphereMetric}
  \source{petersen:page-0135,petersen:page-0136}
  \leanok
  For $a,b>0$, $S^n_a\times S^m_b:=S^n(1/\sqrt a)\times S^m(1/\sqrt b)$, with metric
  $g=\tfrac1a ds_n^2+\tfrac1b ds_m^2$ on $S^n\times S^m$.
\end{definition}

\begin{proposition}[curvature of product spheres]
  \label{prop:pet-ch4-product-sphere-curvature}
  \lean{PetersenLib.productSphereCurvature}
  \source{petersen:page-0136}
  \uses{def:pet-ch4-product-sphere-metric, prop:pet-ch4-euclidean-sphere-constant-curvature, prop:pet-ch4-triple-vanishing-diagonalizes-operator}
  On $S^n_a\times S^m_b$, let $Y$ be a unit vector field on $S^n$, $V$ a unit
  vector field on $S^m$, and $X$ a unit vector field tangent to either factor and
  perpendicular to both $Y$ and $V$. Then $\nabla_YX=0$ if $X$ is tangent to
  $S^m$ and $\nabla_YX$ is tangent to $S^n$ if $X$ is tangent to $S^n$ (so
  $\nabla_YX$ can be computed within $S^n$); consequently
  \[
    \mathfrak{R}(X\wedge Y)=aX\wedge Y\ (X\text{ tangent to }S^n),\qquad
    \mathfrak{R}(X\wedge Y)=0\ (X\text{ tangent to }S^m),\qquad
    \mathfrak{R}(U\wedge V)=bU\wedge V.
  \]
\end{proposition}
\begin{proof}
  \uses{prop:pet-ch4-product-sphere-curvature}
  \sketch
  The Koszul formula shows that the product connection preserves each factor
  and has zero mixed derivatives. Hence pure planes have the curvatures of
  their factors, namely $a$ and $b$ by
  \cref{prop:pet-ch4-euclidean-sphere-constant-curvature}, while mixed planes
  are flat. All components involving three different factor-adapted basis
  directions vanish, so \cref{prop:pet-ch4-triple-vanishing-diagonalizes-operator}
  gives the displayed bivector formulas.
\end{proof}

\begin{remark}[Ricci, scalar curvature, and the Einstein condition for product spheres]
  \label{rem:pet-ch4-product-sphere-ricci-einstein}
  \lean{PetersenLib.productSphereRicciEinstein}
  \source{petersen:page-0136}
  \uses{prop:pet-ch4-product-sphere-curvature, prop:pet-ch4-curvature-operator-diagonal-bounds}
  On $S^n_a\times S^m_b$: $\mathrm{Ric}(X)=(n-1)aX$, $\mathrm{Ric}(V)=(m-1)bV$, and
  $\mathrm{scal}=n(n-1)a+m(m-1)b$; all sectional curvatures lie in
  $[0,\max\{a,b\}]$. Consequently $S^n_a\times S^m_b$ always has constant scalar
  curvature, is Einstein exactly when $(n-1)a=(m-1)b$ (forcing $n,m\ge2$ or
  $n=m=1$), and has constant sectional curvature only when $n=m=1$; its curvature
  tensor is always parallel.
\end{remark}

\subsection{Rotationally Symmetric Metrics}

For $g=dr^2+\rho^2(r)ds_{n-1}^2$, radial derivatives of $\rho$ determine all
curvature components and their Ricci and scalar contractions.

\begin{proposition}[Hessian of $r$ for a rotationally symmetric metric]
  \label{prop:pet-ch4-hess-r-rotsym}
  \lean{PetersenLib.hessDistanceRotSym}
  \source{petersen:page-0136,petersen:page-0137}
  On $g=dr^2+g_r$, $g_r=\rho^2(r)ds_{n-1}^2$, on $(a,b)\times S^{n-1}$, since
  $ds_{n-1}^2$ does not depend on $r$,
  \[
    \mathrm{Hess}\,r=\frac{\partial_r\rho}{\rho}\,g_r.
  \]
\end{proposition}
\begin{proof}
  \uses{prop:pet-ch4-hess-r-rotsym}
  \sketch
  As in \cref{prop:pet-ch4-hess-r-euclidean}, only the warping factor changes
  under radial flow. Consequently
  \[
    2\operatorname{Hess}r=L_{\partial_r}g_r
    =2\rho\dot\rho\,ds_{n-1}^2
    =2\frac{\dot\rho}{\rho}g_r.
  \]
\end{proof}

\begin{lemma}[Lie and covariant derivatives of $\mathrm{Hess}\,r$]
  \label{lem:pet-ch4-hess-r-derivatives}
  \lean{PetersenLib.hessDistanceRotSymDerivatives}
  \source{petersen:page-0137,petersen:page-0138}
  \uses{prop:pet-ch4-hess-r-rotsym}
  With $\mathrm{Hess}\,r$ as in \cref{prop:pet-ch4-hess-r-rotsym},
  \[
    L_{\partial_r}\mathrm{Hess}\,r=\frac{\partial_r^2\rho}{\rho}g_r+(\mathrm{Hess}\,r)^2,
    \qquad
    \nabla_{\partial_i}\mathrm{Hess}\,r=\frac{\partial_r^2\rho}{\rho}g_r-(\mathrm{Hess}\,r)^2,
  \]
  and, restricted to $S^{n-1}$, $\mathrm{Hess}\,r=\tfrac{\partial_r\rho}{\rho}g_r$
  while $R(\cdot,\partial_r,\partial_r,\cdot)=-\tfrac{\partial_r^2\rho}{\rho}g_r$.
\end{lemma}
\begin{proof}
  \uses{lem:pet-ch4-hess-r-derivatives}
  \sketch
  Differentiate $\operatorname{Hess}r=(\dot\rho/\rho)g_r$ and use
  $L_{\partial_r}g_r=2(\dot\rho/\rho)g_r$ from
  \cref{prop:pet-ch4-hess-r-rotsym}. This gives
  $L_{\partial_r}\operatorname{Hess}r=(\ddot\rho/\rho)g_r+
  (\operatorname{Hess}r)^2$; converting Lie to covariant differentiation changes
  the sign of the squared term. The radial curvature equation then yields
  $R(\cdot,\partial_r,\partial_r,\cdot)=-(\ddot\rho/\rho)g_r$.
\end{proof}


\begin{proposition}[normal Jacobi field data and vanishing of the mixed curvature]
  \label{prop:pet-ch4-rotsym-radial-jacobi}
  \lean{PetersenLib.rotSymRadialJacobiData}
  \source{petersen:page-0138}
  \uses{lem:pet-ch4-hess-r-derivatives}
  On $(I\times S^{n-1},dr^2+\rho^2(r)ds_{n-1}^2)$,
  \[
    \nabla_X\partial_r=\begin{cases}\dfrac{\partial_r\rho}{\rho}X&X\text{ tangent to }S^{n-1},\\[4pt]0&X=\partial_r,\end{cases}
    \qquad
    R(X,\partial_r)\partial_r=\begin{cases}-\dfrac{\partial_r^2\rho}{\rho}X&X\text{ tangent to }S^{n-1},\\[4pt]0&X=\partial_r,\end{cases}
  \]
  and the mixed curvature vanishes: $\nabla_X\mathrm{II}=0$ for $\mathrm{II}=\mathrm{Hess}\,r$.
\end{proposition}
\begin{proof}
  \uses{prop:pet-ch4-rotsym-radial-jacobi}
  \sketch
  The formulas for $\nabla_X\partial_r$ and $R(X,\partial_r)\partial_r$ follow directly
  from \cref{prop:pet-ch4-hess-r-rotsym} and
  \cref{lem:pet-ch4-hess-r-derivatives} by
  definition of the second fundamental form and radial curvature. For the mixed
  curvature, since $\mathrm{II}=\tfrac{\partial_r\rho}{\rho}g_r$ and the coefficient
  $\partial_r\rho/\rho$ depends only on $r$ (a function constant on each level set
  of $r$), for $X$ tangent to a level set,
  $\nabla_X\mathrm{II}=D_X(\partial_r\rho/\rho)\,g_r+\tfrac{\partial_r\rho}{\rho}\nabla_Xg_r=0+0=0$,
  since $D_X(\partial_r\rho/\rho)=0$ ($X$ is tangent to the level set on which
  $\partial_r\rho/\rho$ is constant) and $\nabla_Xg_r=0$ ($g_r$ is parallel along its
  own level sets by metric compatibility).
\end{proof}

\begin{theorem}[curvature operator of a rotationally symmetric metric]
  \label{thm:pet-ch4-rotsym-curvature-operator}
  \lean{PetersenLib.rotSymCurvatureOperator}
  \source{petersen:page-0138,petersen:page-0139}
  \uses{prop:pet-ch4-rotsym-radial-jacobi, prop:pet-ch4-triple-vanishing-diagonalizes-operator}
  For $g=dr^2+\rho^2(r)ds_{n-1}^2$ and $X,Y$ tangent to $S^{n-1}$,
  \[
    \mathfrak{R}(X\wedge\partial_r)=-\frac{\ddot\rho}{\rho}X\wedge\partial_r,
    \qquad
    \mathfrak{R}(X\wedge Y)=\frac{1-\dot\rho^2}{\rho^2}X\wedge Y,
  \]
  so $\mathfrak{R}$ is diagonalized, with all sectional curvatures between
  $-\ddot\rho/\rho$ and $(1-\dot\rho^2)/\rho^2$.
\end{theorem}
\begin{proof}
  \uses{thm:pet-ch4-rotsym-curvature-operator}
  \sketch
  The radial formula in \cref{prop:pet-ch4-rotsym-radial-jacobi} gives the
  eigenvalue $-\ddot\rho/\rho$. For tangent vectors, the Gauss equation,
  \cref{prop:pet-ch4-euclidean-sphere-constant-curvature}, and
  $\mathrm{II}=(\dot\rho/\rho)g_r$ give
  \[
    g(R(X,Y)V,W)=\frac{1-\dot\rho^2}{\rho^2}
    g_r(X\wedge Y,W\wedge V).
  \]
  The mixed terms vanish, so
  \cref{prop:pet-ch4-triple-vanishing-diagonalizes-operator} diagonalizes
  $\mathfrak R$ in the radial/tangential frame. Apply
  \cref{prop:pet-ch4-curvature-operator-diagonal-bounds} for the bounds.
\end{proof}


\begin{proposition}[Ricci and scalar curvature of a rotationally symmetric metric]
  \label{prop:pet-ch4-rotsym-ricci-scalar}
  \lean{PetersenLib.rotSymRicciScalar}
  \source{petersen:page-0139}
  \uses{thm:pet-ch4-rotsym-curvature-operator}
  For $g=dr^2+\rho^2(r)ds_{n-1}^2$ on $I\times S^{n-1}$, $n\ge2$, and $X$ tangent
  to $S^{n-1}$,
  \[
    \mathrm{Ric}(X)=\Big((n-2)\frac{1-\dot\rho^2}{\rho^2}-\frac{\ddot\rho}{\rho}\Big)X,
    \qquad
    \mathrm{Ric}(\partial_r)=-(n-1)\frac{\ddot\rho}{\rho}\partial_r,
  \]
  \[
    \mathrm{scal}=-2(n-1)\frac{\ddot\rho}{\rho}+(n-1)(n-2)\frac{1-\dot\rho^2}{\rho^2}.
  \]
\end{proposition}
\begin{proof}
  \uses{prop:pet-ch4-rotsym-ricci-scalar}
  \sketch
  Use an orthonormal frame with first vector $\partial_r$. The two eigenvalues
  in \cref{thm:pet-ch4-rotsym-curvature-operator} contribute once in each
  radial plane and $n-2$ times in the tangential planes through a fixed
  tangential vector. Their sums give the two Ricci eigenvalues. Taking the trace
  adds one radial and $n-1$ tangential Ricci eigenvalues, yielding the stated
  scalar curvature.
\end{proof}

\begin{remark}[the two-dimensional case]
  \label{rem:pet-ch4-rotsym-dim2}
  \lean{PetersenLib.rotSymDim2Curvature}
  \source{petersen:page-0139}
  \uses{prop:pet-ch4-rotsym-ricci-scalar}
  When $n=2$ there are no tangential curvatures, so
  $\mathrm{sec}=-\ddot\rho/\rho$; this is a genuine difference between
  $2$-dimensional and higher-dimensional rotationally symmetric metrics.
\end{remark}

\begin{proposition}[constant curvature of $S^n_k$]
  \label{prop:pet-ch4-snk-constant-curvature}
  \lean{PetersenLib.snkConstantCurvature}
  \source{petersen:page-0139,petersen:page-0140}
  \uses{thm:pet-ch4-rotsym-curvature-operator}
  The metric $dr^2+\mathrm{sn}_k^2(r)ds_{n-1}^2$ has all sectional curvatures
  equal to $k$.
\end{proposition}
\begin{proof}
  \uses{prop:pet-ch4-snk-constant-curvature}
  \sketch
  For $\rho=\operatorname{sn}_k$, the identities
  $\ddot\rho=-k\rho$ and
  $\dot\rho^2=\operatorname{cs}_k^2=1-k\rho^2$ turn both eigenvalues in
  \cref{thm:pet-ch4-rotsym-curvature-operator} into $k$.
\end{proof}

\begin{proposition}[Ricci flat rotationally symmetric metrics are flat]
  \label{prop:pet-ch4-ricci-flat-rotsym}
  \lean{PetersenLib.ricciFlatRotSymIsFlat}
  \source{petersen:page-0139,petersen:page-0140}
  \uses{prop:pet-ch4-rotsym-ricci-scalar}
  If $dr^2+\rho^2(r)ds_{n-1}^2$ is Ricci flat, then the metric is flat. More
  precisely, if $n>2$, then $\rho(r)=a+r$ after reversing $r$ if necessary; if
  $n=2$, then $\rho$ is affine.
\end{proposition}
\begin{proof}
  \uses{prop:pet-ch4-ricci-flat-rotsym}
  \sketch
  The radial Ricci eigenvalue in \cref{prop:pet-ch4-rotsym-ricci-scalar}
  forces $\ddot\rho=0$. For $n>2$, the tangential eigenvalue also forces
  $\dot\rho^2=1$, so reversing $r$ if needed gives $\rho=a+r$. For $n=2$ there
  is no tangential equation, and $\rho=a+br$. In both cases
  $-\ddot\rho/\rho=0$, so every sectional curvature vanishes.
\end{proof}

\begin{definition}[the scalar-flat ODE and its first-order reduction]
  \label{def:pet-ch4-scalar-flat-ode}
  \lean{PetersenLib.scalarFlatRotSymODE}
  \source{petersen:page-0140}
  \uses{prop:pet-ch4-rotsym-ricci-scalar}
  \leanok
  For $n\ge3$, scalar flatness of $dr^2+\rho^2(r)ds_{n-1}^2$ is the equation
  $\ddot\rho+\tfrac{n-2}2\tfrac{\dot\rho^2-1}\rho=0$. Substituting $\dot\rho=G(\rho)$
  (so $\ddot\rho=G'G$) turns this autonomous second-order equation into the
  first-order equation $G'G+\tfrac{n-2}2\tfrac{G^2-1}\rho=0$, whose separated
  solutions satisfy $\dot\rho^2=G^2=1+C\rho^{2-n}$ for a constant $C$, equivalently
  $\ddot\rho=-\tfrac{n-2}2C\rho^{1-n}$ (compatible initial data:
  $\dot\rho(0)^2=1+C\rho(0)^{2-n}$).
\end{definition}

\begin{proposition}[classification of scalar-flat rotationally symmetric solutions]
  \label{prop:pet-ch4-scalar-flat-cases}
  \lean{PetersenLib.scalarFlatRotSymCases}
  \source{petersen:page-0140,petersen:page-0141}
  \uses{def:pet-ch4-scalar-flat-ode}
  For the family $\dot\rho^2=1+C\rho^{2-n}$ of \cref{def:pet-ch4-scalar-flat-ode}:
  if $C=0$ then $\rho(r)=a+r$ (the Euclidean metric); if $C>0$ then $\rho$ is
  concave and its maximal positive interval has a finite endpoint where $\rho$
  reaches $0$ and $\ddot\rho$ becomes singular, so it gives no globally defined
  smooth positive metric; if $C<0$, writing $C=-\rho_0^{n-2}$, the solution with
  $\rho(0)=\rho_0>0$ is even, convex, satisfies $\rho_0\le\rho$ and $|\dot\rho|\le1$,
  and is defined on all of $\mathbb{R}$.
\end{proposition}
\begin{proof}
  \uses{prop:pet-ch4-scalar-flat-cases}
  \sketch
  For $C=0$, the first integral gives $\dot\rho^2=1$ and hence the affine
  Euclidean solution. If $C>0$, then
  $\ddot\rho=-\tfrac{n-2}{2}C\rho^{1-n}<0$; the positive solution reaches
  $\rho=0$ in finite parameter, where the differential equation is singular.
  If $C=-\rho_0^{n-2}$, then
  \[
    \dot\rho^2=1-(\rho_0/\rho)^{n-2},
  \]
  so $\rho\ge\rho_0$, $|\dot\rho|\le1$, and $\ddot\rho>0$. Initial data
  $\rho(0)=\rho_0$, $\dot\rho(0)=0$ give an even solution; the derivative bound
  prevents finite-time escape and yields existence on all of $\mathbb R$.
\end{proof}

\begin{example}[scalar-flat metric on the tautological line bundle over $\mathbb{RP}^{n-1}$]
  \label{ex:pet-ch4-scalar-flat-tautological-bundle}
  \lean{PetersenLib.scalarFlatTautologicalBundleMetric}
  \source{petersen:page-0141}
  \uses{prop:pet-ch4-scalar-flat-cases}
  For the even solution $\rho$ with $\rho(0)=\rho_0>0$ from
  \cref{prop:pet-ch4-scalar-flat-cases}, the map $(r,x)\mapsto(-r,-x)$ is an
  isometry of $\big(\mathbb{R}\times S^{n-1},dr^2+\rho^2(r)ds_{n-1}^2\big)$, giving a
  Riemannian covering $\mathbb{R}\times S^{n-1}\to\tau(\mathbb{RP}^{n-1})$ onto the
  tautological line bundle over $\mathbb{RP}^{n-1}$, and hence a scalar-flat metric on
  $\tau(\mathbb{RP}^{n-1})$. Using $\rho$ as parameter for $r>0$, since
  $d\rho^2=\dot\rho^2dr^2=\big(1-(\rho_0/\rho)^{n-2}\big)dr^2$,
  \[
    dr^2+\rho^2(r)ds_{n-1}^2=\frac1{1-(\rho_0/\rho)^{n-2}}\,d\rho^2+\rho^2ds_{n-1}^2,
  \]
  which approaches the Euclidean metric $d\rho^2+\rho^2ds_{n-1}^2$ as
  $\rho\to\infty$.
\end{example}

\begin{remark}[scalar flat without Ricci flat or constant curvature]
  \label{rem:pet-ch4-scalar-flat-not-ricci-flat}
  \lean{PetersenLib.scalarFlatNotRicciFlatRemark}
  \source{petersen:page-0141}
  \uses{ex:pet-ch4-scalar-flat-tautological-bundle}
  $\mathbb{R}\times S^{n-1}$, $n\ge3$, admits no complete constant-curvature metric
  (later, §5.6.2); moreover if it carried a Ricci-flat metric then $S^{n-1}$
  would carry a Ricci-flat metric too (§7.3.1, Thm.\ 7.3.5), impossible for
  $n=3,4$ (giving flat metrics on $S^2,S^3$, excluded in §5.6.2). Thus the
  tautological bundle of \cref{ex:pet-ch4-scalar-flat-tautological-bundle}
  carries a scalar-flat metric but no Ricci-flat or constant-curvature metric.
\end{remark}

\subsection{Doubly Warped Products}

With two warping functions, the same level-set calculation produces five curvature
eigenvalues, including the mixed-factor term.

\begin{theorem}[curvature operator of a doubly warped product]
  \label{thm:pet-ch4-doubly-warped-curvature}
  \lean{PetersenLib.doublyWarpedCurvatureOperator}
  \source{petersen:page-0142}
  \uses{thm:pet-ch4-rotsym-curvature-operator, prop:pet-ch4-product-sphere-curvature, prop:pet-ch4-triple-vanishing-diagonalizes-operator}
  On $\big(I\times S^p\times S^q,\,dr^2+\rho^2(r)ds_p^2+\phi^2(r)ds_q^2\big)$, the
  second fundamental form of the level sets is
  $\mathrm{II}=(\partial_r\rho)\rho\,ds_p^2+(\partial_r\phi)\phi\,ds_q^2$, the mixed
  curvature vanishes ($\nabla_X\mathrm{II}=0$), and for $X,Y$ tangent to $S^p$,
  $V,W$ tangent to $S^q$,
  \[
    \mathfrak{R}(\partial_r\wedge X)=-\frac{\ddot\rho}\rho\partial_r\wedge X,\quad
    \mathfrak{R}(\partial_r\wedge V)=-\frac{\ddot\phi}\phi\partial_r\wedge V,\quad
    \mathfrak{R}(X\wedge Y)=\frac{1-\dot\rho^2}{\rho^2}X\wedge Y,
  \]
  \[
    \mathfrak{R}(V\wedge W)=\frac{1-\dot\phi^2}{\phi^2}V\wedge W,\qquad
    \mathfrak{R}(X\wedge V)=-\frac{\dot\rho\dot\phi}{\rho\phi}X\wedge V.
  \]
  In particular, all sectional curvatures are convex combinations of
  $-\ddot\rho/\rho$, $-\ddot\phi/\phi$, $(1-\dot\rho^2)/\rho^2$,
  $(1-\dot\phi^2)/\phi^2$, $-\dot\rho\dot\phi/(\rho\phi)$.
\end{theorem}
\begin{proof}
  \uses{thm:pet-ch4-doubly-warped-curvature}
  \sketch
  Applying \cref{prop:pet-ch4-hess-r-rotsym} to both factors gives
  $\mathrm{II}=\rho\dot\rho\,ds_p^2+\phi\dot\phi\,ds_q^2$; as in
  \cref{prop:pet-ch4-rotsym-radial-jacobi}, its levelwise covariant derivative
  vanishes. The radial equation supplies $-\ddot\rho/\rho$ and
  $-\ddot\phi/\phi$. The Gauss equation, together with the round-sphere
  curvature from \cref{prop:pet-ch4-euclidean-sphere-constant-curvature} and
  the single-factor computation \cref{thm:pet-ch4-rotsym-curvature-operator},
  supplies the two pure tangential eigenvalues. The product splitting in
  \cref{prop:pet-ch4-product-sphere-curvature} leaves only the second-fundamental-form
  term on mixed planes, giving $-\dot\rho\dot\phi/(\rho\phi)$. Now apply
  \cref{prop:pet-ch4-triple-vanishing-diagonalizes-operator} and then
  \cref{prop:pet-ch4-curvature-operator-diagonal-bounds}.
\end{proof}

\begin{proposition}[Ricci curvature of a doubly warped product]
  \label{prop:pet-ch4-doubly-warped-ricci}
  \lean{PetersenLib.doublyWarpedRicci}
  \source{petersen:page-0142,petersen:page-0143}
  \uses{thm:pet-ch4-doubly-warped-curvature}
  On $\big(I\times S^p\times S^q,\,dr^2+\rho^2(r)ds_p^2+\phi^2(r)ds_q^2\big)$, for
  $X$ tangent to $S^p$ and $V$ tangent to $S^q$,
  \[
    \mathrm{Ric}(\partial_r)=\Big(-p\frac{\ddot\rho}\rho-q\frac{\ddot\phi}\phi\Big)\partial_r,
  \]
  \[
    \mathrm{Ric}(X)=\Big(-\frac{\ddot\rho}\rho+(p-1)\frac{1-\dot\rho^2}{\rho^2}-q\frac{\dot\rho\dot\phi}{\rho\phi}\Big)X,
    \qquad
    \mathrm{Ric}(V)=\Big(-\frac{\ddot\phi}\phi+(q-1)\frac{1-\dot\phi^2}{\phi^2}-p\frac{\dot\rho\dot\phi}{\rho\phi}\Big)V.
  \]
\end{proposition}
\begin{proof}
  \uses{prop:pet-ch4-doubly-warped-ricci}
  \sketch
  Trace the five curvature eigenvalues of
  \cref{thm:pet-ch4-doubly-warped-curvature} in a frame consisting of
  $\partial_r$, $p$ vectors tangent to $S^p$, and $q$ vectors tangent to $S^q$.
  The radial trace has $p$ copies of $-\ddot\rho/\rho$ and $q$ copies of
  $-\ddot\phi/\phi$; fixing a vector in either sphere factor gives one radial,
  the remaining pure-factor, and all mixed-factor contributions displayed in
  the statement. This is the same contraction as
  \cref{prop:pet-ch4-rotsym-ricci-scalar}.
\end{proof}

\subsection{The Schwarzschild Metric}

The Ricci-flat equations for a doubly warped metric with a circle factor reduce to
the scalar-flat radial equation.

\begin{definition}[the Schwarzschild Ricci-flat equations]
  \label{def:pet-ch4-schwarzschild-construction}
  \lean{PetersenLib.schwarzschildRicciFlatEquations}
  \source{petersen:page-0143}
  \uses{prop:pet-ch4-doubly-warped-ricci}
  \leanok
  For $n\ge4$, take $p=n-2$, $q=1$ in
  \cref{prop:pet-ch4-doubly-warped-ricci} on
  $(0,\infty)\times S^{n-2}\times S^1$: Ricci flatness $\mathrm{Ric}\equiv0$
  ($\mathrm{Ric}(\partial_r)=0$, $\mathrm{Ric}(X)=0$, $\mathrm{Ric}(V)=0$) becomes
  \[
    -(n-2)\frac{\ddot\rho}\rho-\frac{\ddot\phi}\phi=0,\qquad
    -\frac{\ddot\rho}\rho+(n-3)\frac{1-\dot\rho^2}{\rho^2}-\frac{\dot\rho\dot\phi}{\rho\phi}=0,\qquad
    -\frac{\ddot\phi}\phi-(n-2)\frac{\dot\rho\dot\phi}{\rho\phi}=0.
  \]
\end{definition}

\begin{proposition}[reduction to the scalar-flat equation]
  \label{prop:pet-ch4-schwarzschild-reduces-scalar-flat}
  \lean{PetersenLib.schwarzschildReducesToScalarFlat}
  \source{petersen:page-0143}
  \uses{def:pet-ch4-schwarzschild-construction, def:pet-ch4-scalar-flat-ode}
  \leanok
  For $n\ge4$, subtracting the first and third equations of
  \cref{def:pet-ch4-schwarzschild-construction} gives
  $\ddot\rho/\rho=\dot\rho\dot\phi/(\rho\phi)$, which integrates to
  $\dot\rho=c\phi$ for a constant $c$; substituting
  $\dot\rho\dot\phi/(\rho\phi)=\ddot\rho/\rho$ back into the second equation
  reduces it to the scalar-flat equation
  $-2\ddot\rho/\rho+(n-3)(1-\dot\rho^2)/\rho^2=0$ for $dr^2+\rho^2(r)ds_{n-2}^2$
  (the case of \cref{def:pet-ch4-scalar-flat-ode} with $n$ replaced by $n-1$).
\end{proposition}
\begin{proof}
  \uses{prop:pet-ch4-schwarzschild-reduces-scalar-flat}
  \leanok
  \sketch
  Subtracting the third equation of
  \cref{def:pet-ch4-schwarzschild-construction} from the first gives
  $\ddot\rho/\rho=\dot\rho\dot\phi/(\rho\phi)$. Hence
  $d(\log|\dot\rho|)/dr=d(\log|\phi|)/dr$, so $\dot\rho=c\phi$. Replacing
  $\dot\rho\dot\phi/(\rho\phi)$ by $\ddot\rho/\rho$ in the second equation
  yields
  $-2\ddot\rho/\rho+(n-3)(1-\dot\rho^2)/\rho^2=0$, namely
  \cref{def:pet-ch4-scalar-flat-ode} in dimension $n-1$.
\end{proof}


\begin{proposition}[the Schwarzschild metric is Ricci flat]
  \label{prop:pet-ch4-schwarzschild-ricci-flat}
  \lean{PetersenLib.schwarzschildRicciFlat}
  \source{petersen:page-0143,petersen:page-0144}
  \uses{prop:pet-ch4-schwarzschild-reduces-scalar-flat, prop:pet-ch4-scalar-flat-cases}
  For $n\ge4$, let $\rho$ be the even solution of the scalar-flat equation from
  \cref{prop:pet-ch4-scalar-flat-cases} (with $n$ there replaced by $n-1$) with
  $\rho(0)=\rho_0>0$ and $\dot\rho^2=1-(\rho_0/\rho)^{n-3}$, and define $\phi$ by
  $\dot\rho=c\phi$ with $c=\tfrac{n-3}2\rho_0^{-1}$. Then $\phi$ is odd,
  $\phi(0)=0$, $\dot\phi(0)=1$, $\dot\phi=(\rho_0/\rho)^{n-2}$, and $(\rho,\phi)$
  solve the equations of \cref{def:pet-ch4-schwarzschild-construction}; the
  resulting metric $dr^2+\rho^2(r)ds_{n-2}^2+\phi^2(r)d\theta^2$ is a Ricci-flat
  metric on $\mathbb{R}^2\times S^{n-2}$ (the \emph{Schwarzschild metric}), with the
  explicit algebraic form
  \[
    \frac1{1-(\rho_0/\rho)^{n-3}}\,d\rho^2+\rho^2ds_{n-2}^2
    +\rho_0^2\frac4{(n-3)^2}\Big(1-\Big(\frac{\rho_0}\rho\Big)^{n-3}\Big)\,d\theta^2,
  \]
  which is asymptotic to the flat metric on $\mathbb{R}^{n-1}\times S^1$ (suitably
  scaled) as $\rho\to\infty$.
\end{proposition}
\begin{proof}
  \uses{prop:pet-ch4-schwarzschild-ricci-flat}
  \sketch
  Differentiating $\dot\rho^2=1-\rho_0^{n-3}\rho^{3-n}$ gives

  \[
    2\ddot\rho=(n-3)\rho_0^{n-3}\rho^{2-n}.
  \]
  Thus $\dot\rho=c\phi$ and $\dot\phi(0)=1$ force
  $c=\tfrac{n-3}{2\rho_0}$ and
  $\dot\phi=(\rho_0/\rho)^{n-2}$. Since $\rho$ is even, $\phi$ is odd.
  A further derivative yields
  $\ddot\phi=-(n-2)\dot\rho\dot\phi/\rho$, while
  $\ddot\rho/\rho=\dot\rho\dot\phi/(\rho\phi)$. These identities verify the
  first and third equations of \cref{def:pet-ch4-schwarzschild-construction};
  the second is the scalar-flat reduction in
  \cref{prop:pet-ch4-schwarzschild-reduces-scalar-flat}. Finally,
  $d\rho^2=\dot\rho^2dr^2$ and
  $\phi^2=4\rho_0^2(1-(\rho_0/\rho)^{n-3})/(n-3)^2$ give the stated algebraic
  form, as in \cref{ex:pet-ch4-scalar-flat-tautological-bundle}; its coefficients
  have the asserted flat limits as $\rho\to\infty$.
\end{proof}


\begin{remark}[the classical space-time Schwarzschild metric]
  \label{rem:pet-ch4-schwarzschild-physical}
  \lean{PetersenLib.schwarzschildPhysicalMetric}
  \source{petersen:page-0144}
  \uses{prop:pet-ch4-schwarzschild-ricci-flat}
  Replacing $S^1$ by $\mathbb{R}$, treating the parameter $c$ (the speed of light) as
  independent of $\rho_0$, and passing to a Lorentzian signature gives the
  classical (non-Riemannian) Schwarzschild space-time metric
  \[
    \frac1{1-\rho_0/\rho}\,d\rho^2+\rho^2ds_2^2-\frac1{c^2}\Big(1-\frac{\rho_0}\rho\Big)dt^2,
  \]
  which is not smooth at $\rho=\rho_0$.
\end{remark}

\section{Warped Products in General}

Warped products admit a potential-function description, local Hessian
characterizations, and conformal models.

\subsection{Basic Constructions}

\begin{definition}[warped product over an interval]
  \label{def:pet-ch4-warped-product-general}
  \lean{PetersenLib.warpedProductGeneral}
  \source{petersen:page-0145}
  \leanok
  Given a Riemannian manifold $(H,g_H)$ and an open interval $I\subset\mathbb{R}$, a
  \emph{warped product} (over $I$) is $g=dr^2+\rho^2(r)g_H$ on $I\times H$, with
  $\rho>0$ on $I$. The more general form $\psi^2(r)dr^2+\rho^2(r)g_H$ reduces to
  this form by the change of variable $d\rho'=\psi(r)dr$, giving
  $d(\rho')^2+\rho^2(r(\rho'))g_H$. Basic special cases: the plain product
  $g=dr^2+g_H$, and Euclidean polar coordinates $dr^2+r^2ds_{n-1}^2$ on
  $(0,\infty)\times S^{n-1}$.
\end{definition}

\begin{definition}[the potential function $f$ of a warped product]
  \label{def:pet-ch4-warping-potential-f}
  \lean{PetersenLib.warpedProductPotential}
  \source{petersen:page-0145}
  \uses{def:pet-ch4-warped-product-general}
  \leanok
  For $g=dr^2+\rho^2(r)g_H$ on $M=I\times H$, set $f=\int\rho\,dr$, so $df=\rho\,dr$
  and
  \[
    g=dr^2+\rho^2(r)g_H=\frac1{\rho^2(r)}df^2+\rho^2(r)g_H.
  \]
\end{definition}

\begin{proposition}[the Hessian of the potential is conformal to $g$]
  \label{prop:pet-ch4-hess-f-conformal}
  \lean{PetersenLib.hessPotentialConformal}
  \source{petersen:page-0145,petersen:page-0146}
  \uses{def:pet-ch4-warping-potential-f, prop:pet-ch4-hess-r-rotsym}
  \dcref{ch4:4.3.1}
  For $f=\int\rho\,dr$ as in \cref{def:pet-ch4-warping-potential-f},
  $\mathrm{Hess}\,f=\dot\rho\,g$.
\end{proposition}
\begin{proof}
  \uses{prop:pet-ch4-hess-f-conformal}
  \sketch
  Since $df=\rho\,dr$,
  \[
    \operatorname{Hess}f=\dot\rho\,dr^2+\rho\operatorname{Hess}r.
  \]
  Substitute \cref{prop:pet-ch4-hess-r-rotsym}, namely
  $\operatorname{Hess}r=\dot\rho\rho\,g_H$, to obtain
  $\operatorname{Hess}f=\dot\rho(dr^2+\rho^2g_H)=\dot\rho\,g$.
\end{proof}

\begin{remark}[warped product data from $f$ and $|\nabla f|$ alone]
  \label{rem:pet-ch4-warped-product-canonical-form}
  \lean{PetersenLib.warpedProductFromPotentialAndGradientNorm}
  \source{petersen:page-0146}
  \uses{prop:pet-ch4-hess-f-conformal}
  With $\rho=|\nabla f|$ (so $\rho^2=g(\nabla f,\nabla f)$),
  $\dot\rho=\tfrac{d\rho}{dr}=\tfrac{d\rho}{df}\tfrac{df}{dr}=\tfrac{d\rho}{df}\rho
  =\tfrac12\tfrac{d|\nabla f|^2}{df}$; hence the warped product representation
  depends only on $f$ and $|\nabla f|$:
  \[
    g=\frac1{|\nabla f|^2}df^2+|\nabla f|^2g_H,\qquad
    \mathrm{Hess}\,f=\frac12\frac{d|\nabla f|^2}{df}\,g.
  \]
\end{remark}

\begin{example}[the potential for constant-curvature warped products]
  \label{ex:pet-ch4-hess-f-space-forms}
  \lean{PetersenLib.hessFPotentialSpaceForms}
  \source{petersen:page-0146,petersen:page-0147}
  \uses{prop:pet-ch4-hess-f-conformal, prop:pet-ch4-snk-constant-curvature}
  \dcref{ch4:4.3.2}
  For $dr^2+\mathrm{sn}_k^2(r)ds_{n-1}^2$, with $f$ the antiderivative of
  $\mathrm{sn}_k$ vanishing at $r=0$: when $k=0$, $f=\tfrac12r^2$ and
  $\mathrm{Hess}\,f=g$; when $k\ne0$, $f=\tfrac1k-\tfrac1k\mathrm{cs}_k(r)$ and
  $\mathrm{Hess}\,f=\mathrm{cs}_k(r)\,g=(1-kf)\,g$. In particular, $k=1$ gives
  $f=1-\cos r$, $\mathrm{Hess}\,f=(1-f)g$, and $k=-1$ gives $f=\cosh r-1$,
  $\mathrm{Hess}\,f=(1+f)g$.
\end{example}

\subsection{General Characterization}

The conformal-Hessian condition forces a warped product away from critical points;
the nondegenerate critical-point case fixes the round angular metric.

\begin{theorem}[conformal Hessian gives a warped product away from critical points]
  \label{thm:pet-ch4-brinkmann-warped-product}
  \lean{PetersenLib.brinkmannWarpedProductAwayFromCriticalPoints}
  \source{petersen:page-0147}
  \uses{def:pet-ch4-warped-product-general}
  \dcref{ch4:4.3.3a}
  If $f$ is a smooth function on $(M,g)$ with $\mathrm{Hess}\,f=\lambda g$ for a
  function $\lambda$, then around any point $p$ with $df|_p\ne0$, $(M,g)$ is
  locally a warped product $g=dr^2+\rho^2(r)g_H$.
\end{theorem}
\begin{proof}
  \uses{thm:pet-ch4-brinkmann-warped-product}
  \sketch
  Put $\rho=|\nabla f|$. From $\operatorname{Hess}f=\lambda g$,
  $d(\rho^2)=2\lambda\,df$, so $\rho$ and $\lambda$ are constant on regular
  level sets of $f$. The closed form $\rho^{-1}df$ defines a coordinate $r$
  with $|dr|=1$, and the metric splits as $g=dr^2+g_r$. On a level set,
  \[
    \operatorname{Hess}r=\frac{\lambda}{\rho}g_r,
    \qquad L_{\partial_r}g_r=2\frac{\lambda}{\rho}g_r.
  \]
  Since $\partial_r\rho=\lambda$, this evolution equation integrates to
  $g_r=\rho^2(r)g_H$, which is the required local warped product.
\end{proof}


\begin{theorem}[conformal Hessian at a nondegenerate critical point]
  \label{thm:pet-ch4-brinkmann-critical-point}
  \lean{PetersenLib.brinkmannAtNondegenerateCriticalPoint}
  \source{petersen:page-0147,petersen:page-0148,petersen:page-0149}
  \uses{thm:pet-ch4-brinkmann-warped-product}
  \dcref{ch4:4.3.3b}
  If $\mathrm{Hess}\,f=\lambda g$, $df|_p=0$, and $\lambda(p)\ne0$, then after
  replacing $f$ by an affine rescaling one may assume $f(p)=0$, $\lambda(p)=1$;
  if $M$ is (locally) the connected component of $\{f<\varepsilon\}$ containing
  $p$ with $p$ the only critical point, then near $p$,
  $g=dr^2+\rho^2(r)ds_{n-1}^2$ for the standard round metric $ds_{n-1}^2$ on
  $S^{n-1}$, with $\rho(0)=0$, $\dot\rho(0)=1$.
\end{theorem}
\begin{proof}
  \uses{thm:pet-ch4-brinkmann-critical-point}
  \sketch
  The Morse lemma gives coordinates centered at $p$ with
  $f=\tfrac12\sum_i(y^i)^2$, so nearby regular level sets are spheres. Applying
  \cref{thm:pet-ch4-brinkmann-warped-product} off $p$ gives
  $g=dr^2+\rho^2g_{n-1}$ and
  $f'=\rho$, $\rho'=\lambda$, with
  $\rho(0)=0$ and $\rho'(0)=1$. Smoothness at the collapsed level determines
  the angular metric: restricting to every radial two-plane and writing
  $g_{S^1}=\phi^2(\theta)d\theta^2$, the Cartesian metric components have a
  direction-independent limit only when $\phi=1$. Repeating this along every
  great circle gives $g_{n-1}=ds_{n-1}^2$.
\end{proof}


\begin{corollary}[characterization of warped-product constant-curvature metrics]
  \label{cor:pet-ch4-brinkmann-constant-curvature}
  \lean{PetersenLib.brinkmannConstantCurvatureCorollary}
  \source{petersen:page-0149,petersen:page-0150}
  \uses{thm:pet-ch4-brinkmann-critical-point, ex:pet-ch4-hess-f-space-forms}
  \dcref{ch4:4.3.4}
  If $f(p)=0$, $df|_p=0$, and $\mathrm{Hess}\,f=(1-kf)g$, then near $p$ the metric
  is the warped-product metric of curvature $k$ from
  \cref{ex:pet-ch4-hess-f-space-forms}.
\end{corollary}
\begin{proof}
  \uses{cor:pet-ch4-brinkmann-constant-curvature}
  \sketch
  By \cref{thm:pet-ch4-brinkmann-critical-point}, the metric has round angular
  part and warping function $\rho=f'$. The initial-value problem
  $f''=1-kf$, $f(0)=f'(0)=0$ has the solutions listed in
  \cref{ex:pet-ch4-hess-f-space-forms}: $\rho=r$ for $k=0$ and
  $\rho=\operatorname{sn}_k(r)$ otherwise. These are precisely the local
  constant-curvature models.
\end{proof}

\begin{remark}[transnormal functions]
  \label{rem:pet-ch4-transnormal-functions}
  \lean{PetersenLib.transnormalFunctionRemark}
  \source{petersen:page-0150}
  \uses{thm:pet-ch4-brinkmann-warped-product}
  \dcref{ch4:4.3.5}
  A function $f:M\to\mathbb{R}$ is \emph{transnormal} if $|df|^2=\rho^2(f)$ for some
  smooth $\rho$; functions with conformal Hessian are locally transnormal by
  \cref{thm:pet-ch4-brinkmann-warped-product}, but the converse fails: on the
  doubly warped representation $dr^2+\sin^2(r)d\theta_1^2+\cos^2(r)d\theta_2^2$
  of $S^3(1)$ on $(0,\pi/2)\times S^1\times S^1$, $f=\tfrac14\sin(2r)$ is
  transnormal but does not have conformal Hessian.
\end{remark}

\subsection{Conformal Representations of Warped Products}

Two changes of radial variable express a warped product as a conformal multiple of
either a product metric or Euclidean polar coordinates.

\begin{definition}[conformal change of metric]
  \label{def:pet-ch4-conformal-change}
  \lean{PetersenLib.conformalChangeOfMetric}
  \source{petersen:page-0150}
  \leanok
  For a Riemannian manifold $(M,g)$ and positive $\psi$ on $M$, $(M,\psi^2g)$ is
  a \emph{conformal change} of $(M,g)$, with $\psi^2$ the \emph{conformal
  factor}.
\end{definition}

\begin{definition}[the two conformal forms of a warped product]
  \label{def:pet-ch4-warped-product-conformal-forms}
  \lean{PetersenLib.warpedProductConformalForms}
  \source{petersen:page-0150,petersen:page-0151}
  \uses{def:pet-ch4-conformal-change, def:pet-ch4-warped-product-general}
  A warped product $dr^2+\rho^2(r)g_H$ can be written conformally in two ways,
  via $dr=\psi(\rho)d\rho$:
  \[
    dr^2+\rho^2(r)g_H=\psi^2(\rho)\big(d\rho^2+g_H\big)\quad\text{when }\rho(r)=\psi(\rho),
  \]
  \[
    dr^2+\rho^2(r)g_H=\psi^2(\rho)\big(d\rho^2+\rho^2g_H\big)\quad\text{when }\rho(r)=\rho\psi(\rho).
  \]
\end{definition}

\paragraph{4.3.3.1 Conformal Models of Spheres}

\begin{proposition}[the Mercator conformal model of the sphere]
  \label{prop:pet-ch4-mercator-model}
  \lean{PetersenLib.mercatorConformalModel}
  \source{petersen:page-0151,petersen:page-0152}
  \uses{def:pet-ch4-warped-product-conformal-forms, prop:pet-ch4-snk-constant-curvature}
  Writing $R^2ds_n^2=dr^2+R^2\sin^2(r/R)ds_{n-1}^2$ (curvature $1/R^2$) in the
  first conformal form of \cref{def:pet-ch4-warped-product-conformal-forms},
  with $\psi(\rho)d\rho=dr$ and $\psi(\rho)=R\sin(r/R)$,
  \[
    \rho=\frac12\log\frac{1-\cos(r/R)}{1+\cos(r/R)},\qquad
    \cos(r/R)=\frac{1-e^{2\rho}}{1+e^{2\rho}},\qquad
    \psi^2(\rho)=\frac{4R^2e^{2\rho}}{(1+e^{2\rho})^2},
  \]
  so $R^2ds_n^2=\dfrac{4R^2e^{2\rho}}{(1+e^{2\rho})^2}\big(d\rho^2+ds_{n-1}^2\big)$.
\end{proposition}
\begin{proof}
  \uses{prop:pet-ch4-mercator-model}
  \sketch
  The first form in \cref{def:pet-ch4-warped-product-conformal-forms} requires
  $d\rho=dr/(R\sin(r/R))$. Integration gives
  $e^{2\rho}=(1-\cos(r/R))/(1+\cos(r/R))$. Solving for the cosine and using
  $\sin^2=1-\cos^2$ yields
  $\cos(r/R)=(1-e^{2\rho})/(1+e^{2\rho})$ and
  $\psi^2=4R^2e^{2\rho}/(1+e^{2\rho})^2$.
\end{proof}

\begin{proposition}[the polar-coordinate (stereographic) conformal model of the sphere]
  \label{prop:pet-ch4-sphere-stereographic-conformal}
  \lean{PetersenLib.sphereStereographicConformalModel}
  \source{petersen:page-0152,petersen:page-0153}
  \uses{def:pet-ch4-warped-product-conformal-forms, prop:pet-ch4-mercator-model}
  In the second conformal form of
  \cref{def:pet-ch4-warped-product-conformal-forms}, applied to
  $R^2ds_n^2=dr^2+R^2\sin^2(r/R)ds_{n-1}^2$,
  \[
    \cos(r/R)=\frac{1-\rho^2}{1+\rho^2},\qquad
    \psi^2(\rho)=\frac{4R^2}{(1+\rho^2)^2},
  \]
  so
  \[
    R^2ds_n^2=\frac{4R^2}{(1+\rho^2)^2}\big(d\rho^2+\rho^2ds_{n-1}^2\big)
    =\frac{4R^2}{(1+\rho^2)^2}g_{\mathrm{eu}^n}.
  \]
\end{proposition}
\begin{proof}
  \uses{prop:pet-ch4-sphere-stereographic-conformal}
  \sketch
  In the second form of
  \cref{def:pet-ch4-warped-product-conformal-forms}, integration of
  $d\rho/\rho=dr/(R\sin(r/R))$ gives
  $\rho^2=(1-\cos(r/R))/(1+\cos(r/R))$. Hence
  $\cos(r/R)=(1-\rho^2)/(1+\rho^2)$ and
  $R^2\sin^2(r/R)=4R^2\rho^2/(1+\rho^2)^2$. Dividing by $\rho^2$ gives the
  conformal factor; the remaining factor is the Euclidean polar metric. The
  algebra is the radial version of \cref{prop:pet-ch4-mercator-model}.
\end{proof}

\paragraph{4.3.3.2 Conformal Models of Hyperbolic Space}

\begin{proposition}[the Poincaré ball model]
  \label{prop:pet-ch4-poincare-disc-model}
  \lean{PetersenLib.poincareDiscModel}
  \source{petersen:page-0153}
  \uses{def:pet-ch4-warped-product-conformal-forms, prop:pet-ch4-sphere-stereographic-conformal}
  Writing hyperbolic space as
  $dr^2+\mathrm{sn}_{-1/R^2}^2(r)ds_{n-1}^2=dr^2+R^2\sinh^2(r/R)ds_{n-1}^2
  =R^2\big(dt^2+\sinh^2(t)ds_{n-1}^2\big)$, the same construction as
  \cref{prop:pet-ch4-sphere-stereographic-conformal} (with $\sinh,\cosh$ in
  place of $\sin,\cos$) gives the \emph{Poincaré model} on the unit ball,
  \[
    R^2\big(dt^2+\sinh^2(t)ds_{n-1}^2\big)=\frac{4R^2}{(1-\rho^2)^2}\,g_{\mathrm{eu}^n},
    \qquad\rho\in B(0,1).
  \]
\end{proposition}

\begin{definition}[the upper half-space model]
  \label{def:pet-ch4-upper-half-space-model}
  \lean{PetersenLib.upperHalfSpaceModel}
  \source{petersen:page-0153}
  \uses{def:pet-ch4-warped-product-general}
  \leanok
  On the half-space $x^n>0$, the metric $(1/x^n)^2\big((dx^1)^2+\dots+(dx^n)^2\big)$
  becomes, with $r=\log(x^n)$, the warped product
  $g=dr^2+(e^{-r})^2\big((dx^1)^2+\dots+(dx^{n-1})^2\big)$.
\end{definition}

\begin{theorem}[inversion transforms the upper half-space model to the Poincaré ball model]
  \label{thm:pet-ch4-inversion-isometry}
  \lean{PetersenLib.inversionUpperHalfSpaceToPoincareball}
  \source{petersen:page-0153,petersen:page-0154}
  \uses{prop:pet-ch4-poincare-disc-model, def:pet-ch4-upper-half-space-model}
  Write $x=(x^1,\dots,x^{n-1})\in\mathbb{R}^{n-1}$, $y=x^n>0$. The inversion in the
  sphere of radius $\sqrt2$ centered at $(0,-1)\in\mathbb{R}^{n-1}\times\mathbb{R}$,
  \[
    F(x,y)=\frac1{r^2}\big(2x,\,1-|x|^2-y^2\big),\qquad r^2=|x|^2+(y+1)^2,
  \]
  maps the upper half-space $\{y>0\}$ into the unit ball ($|F(x,y)|^2=1-4y/r^2=:\rho^2$)
  and pulls back the Poincaré ball metric $\tfrac4{(1-\rho^2)^2}g_{\mathrm{eu}^n}$
  to the upper half-space metric $\tfrac1{y^2}\big(dy^2+g_{\mathbb{R}^{n-1}}\big)$.
\end{theorem}
\begin{proof}
  \uses{thm:pet-ch4-inversion-isometry}
  \sketch
  Differentiating the displayed inversion gives
  $dF^k=2\,dx^k/r^2-4x^k\,dr/r^3$ for $k<n$ and the analogous expression with
  $y+1$ for $dF^n$, where $r\,dr=x\cdot dx+(y+1)dy$. Substitution into the
  Euclidean norm cancels the $dr$-terms and gives
  \[
    F^*g_{\mathrm{eu}^n}=\frac4{r^4}
    \left(dy^2+\sum_{k<n}(dx^k)^2\right).
  \]
  Since $1-|F|^2=4y/r^2$, multiplying by $4/(1-|F|^2)^2$ yields
  $y^{-2}(dy^2+g_{\mathbb R^{n-1}})$, the metric in
  \cref{def:pet-ch4-upper-half-space-model}.
\end{proof}

\begin{remark}[the Riemann model of constant curvature]
  \label{rem:pet-ch4-riemann-model}
  \lean{PetersenLib.riemannConstantCurvatureModel}
  \source{petersen:page-0155}
  \uses{def:pet-ch4-conformal-change}
  For $\psi^2\big((dx^1)^2+\dots+(dx^n)^2\big)$ to have constant curvature $k$,
  taking $\psi=\big(1+\tfrac k4r^2\big)^{-1}$ (with $\psi\cdot dx^1,\dots,\psi\cdot dx^n$
  the orthonormal coframe and $\psi^{-1}\partial_1,\dots,\psi^{-1}\partial_n$ the
  orthonormal frame used in the Koszul-formula computation of $R$) gives the
  \emph{Riemann model}, defined on all of $\mathbb{R}^n$ if $k\ge0$ and on
  $B(0,2/\sqrt{|k|})$ if $k<0$.
\end{remark}

\begin{proposition}[the Riemann model with $k=-1$ is isometric to the Poincaré model]
  \label{prop:pet-ch4-riemann-poincare-isometric}
  \lean{PetersenLib.riemannPoincareIsometric}
  \source{petersen:page-0155}
  \uses{rem:pet-ch4-riemann-model, prop:pet-ch4-poincare-disc-model}
  The map $F(x)=2x$ carries the Poincaré model on the unit ball to the Riemann
  model with $k=-1$ on the ball of radius $2$:
  \[
    \frac1{\big(1-\tfrac14|F|^2\big)^2}\Big(\sum_{k=1}^n(dF^k)^2\Big)
    =\frac4{(1-|x|^2)^2}\Big(\sum_{k=1}^n(dx^k)^2\Big).
  \]
\end{proposition}
\begin{proof}
  \uses{prop:pet-ch4-riemann-poincare-isometric}
  \sketch
  Under $F(x)=2x$, one has $|F|^2=4|x|^2$ and
  $\sum_k(dF^k)^2=4\sum_k(dx^k)^2$. Substitution in the Riemann conformal factor
  gives exactly $4(1-|x|^2)^{-2}g_{\mathrm{eu}^n}$.
\end{proof}

\subsection{Singular Points}

\begin{proposition}[smoothness criterion in the conformal polar-coordinate model]
  \label{prop:pet-ch4-conformal-smoothness-criterion}
  \lean{PetersenLib.conformalSmoothnessCriterion}
  \source{petersen:page-0155}
  \uses{def:pet-ch4-warped-product-conformal-forms}
  For $dr^2+\varphi^2(r)ds_{n-1}^2=\psi^2(\rho)\big(d\rho^2+\rho^2ds_{n-1}^2\big)$
  with $\varphi(r_0)=0$, $\rho(r_0)=0$ at $r_0\in\partial I$: the right-hand side
  is smooth at $\rho=0$ iff $\psi(0)>0$ and $\psi^{(\mathrm{odd})}(0)=0$; translated
  back to $\varphi$, this is the usual criterion $\dot\varphi(r_0)=\pm1$,
  $\varphi^{(\mathrm{even})}(r_0)=0$.
\end{proposition}
\begin{proof}
  \uses{prop:pet-ch4-conformal-smoothness-criterion}
  \sketch
  A radial conformal factor extends smoothly across the Euclidean origin exactly
  when it is the restriction of a smooth even function of $\rho$, with positive
  value there. This is equivalent to $\psi(0)>0$ and the vanishing of all odd
  derivatives. Translating Taylor series through either radial substitution in
  \cref{def:pet-ch4-warped-product-conformal-forms} gives
  $\dot\varphi(r_0)=\pm1$ and vanishing even derivatives of $\varphi$.
\end{proof}

\section{Metrics on Lie Groups}

Left translations identify a Lie-group metric with an inner product on its Lie
algebra; biinvariance is characterized by the adjoint action.

\subsection{Generalities on Left-invariant Metrics}

\begin{definition}[left-invariant metric via left translation]
  \label{def:pet-ch4-left-invariant-metric-lie}
  \lean{PetersenLib.leftInvariantMetricViaTranslation}
  \source{petersen:page-0155}
  \leanok
  Fixing an inner product $(\cdot,\cdot)$ on $T_eG$ for a Lie group $G$ and
  translating it to $T_gG$ via $L_g(x)=gx$ produces a Riemannian metric $(\cdot,\cdot)$
  on $G$; left-invariant vector fields $X$ (those with $DL_g(X|_h)=X|_{gh}$)
  are determined by their value at $e$, identifying $T_eG$ with the space
  $\mathfrak g$ of left-invariant vector fields, which is a vector space and a
  Lie algebra (bracket of left-invariant fields is left-invariant); on matrix
  groups this bracket is the matrix commutator.
\end{definition}

\begin{proposition}[left translations are isometries]
  \label{prop:pet-ch4-left-translation-isometry}
  \lean{PetersenLib.leftTranslationIsIsometry}
  \source{petersen:page-0155}
  \uses{def:pet-ch4-left-invariant-metric-lie}
  \leanok
  With the metric of \cref{def:pet-ch4-left-invariant-metric-lie}, $L_g$ is a
  Riemannian isometry for every $g\in G$.
\end{proposition}
\begin{proof}
  \uses{prop:pet-ch4-left-translation-isometry}
  \leanok
  \sketch
  The metric at $h$ is the pushforward of the fixed inner product at $e$ by
  $(DL_h)_e$. Since $L_g\circ L_h=L_{gh}$, the differential $(DL_g)_h$ carries
  this pushforward to the one defining the metric at $gh$; it is therefore an
  isometry at every $h$.
\end{proof}

\begin{definition}[the exponential map and its integral-curve property]
  \label{def:pet-ch4-lie-algebra-exp}
  \lean{PetersenLib.lieAlgebraExpMap}
  \source{petersen:page-0155,petersen:page-0156}
  \uses{def:pet-ch4-left-invariant-metric-lie}
  \leanok
  For $X\in\mathfrak g$, $\exp(tX)$ denotes the integral curve of $X$ through
  $e\in G$; on a matrix group the matrix exponential $e^{tX}$ is this integral
  curve, and the flow of $X$ is $F^t(x)=x\exp(tX)=L_x(\exp(tX))=R_{\exp(tX)}(x)$.
\end{definition}

\begin{proposition}[the matrix exponential is the integral curve of a left-invariant field]
  \label{prop:pet-ch4-matrix-exp-integral-curve}
  \lean{PetersenLib.matrixExpIsIntegralCurve}
  \source{petersen:page-0155,petersen:page-0156}
  \uses{def:pet-ch4-lie-algebra-exp}
  \leanok
  On a matrix Lie group, $\tfrac{d}{dt}\big|_{t=t_0}e^{tX}=X|_{e^{t_0X}}$; more
  generally $t\mapsto\exp(tX)$ is the unique integral curve of $X$ through $e$
  because $t\mapsto\exp(tX)$ is a homomorphism $\mathbb{R}\to G$ with derivative
  $X|_e$ at $t=0$.
\end{proposition}
\begin{proof}
  \uses{prop:pet-ch4-matrix-exp-integral-curve}
  \leanok
  \sketch
  From $e^{(t_0+s)X}=L_{e^{t_0X}}(e^{sX})$,
  \[
    \frac d{dt}\Big|_{t=t_0}e^{tX}
    =(DL_{e^{t_0X}})_eX=X|_{e^{t_0X}}.
  \]
  For an arbitrary Lie group, the same conclusion follows because
  $t\mapsto\exp(tX)$ is the unique one-parameter subgroup with initial
  derivative $X$.
\end{proof}

\begin{remark}[right-invariant fields are the Killing fields of a left-invariant metric]
  \label{rem:pet-ch4-flow-right-invariant-killing}
  \lean{PetersenLib.rightInvariantFieldsAreKilling}
  \source{petersen:page-0156}
  \uses{prop:pet-ch4-left-translation-isometry, def:pet-ch4-lie-algebra-exp}
  For $X\in\mathfrak g$, the flow $F^t(x)=x\exp(tX)=R_{\exp(tX)}(x)$ need not
  act by isometries of a merely left-invariant metric, unless the metric is
  also invariant under right translations (\emph{biinvariant}); thus elements
  of $\mathfrak g$ are not in general Killing fields. Instead, it is the
  right-invariant fields -- with flow
  $F^t(x)=\exp(tX)x=L_{\exp(tX)}(x)$ -- that are Killing fields for a
  left-invariant metric, by \cref{prop:pet-ch4-left-translation-isometry}.
\end{remark}

\begin{definition}[the adjoint action of $G$ on $\mathfrak g$]
  \label{def:pet-ch4-adjoint-action}
  \lean{PetersenLib.adjointActionOfLieGroup}
  \source{petersen:page-0156}
  \uses{def:pet-ch4-lie-algebra-exp}
  \leanok
  Conjugation $\mathrm{Ad}_g(x)=gxg^{-1}$ is the \emph{adjoint action of $G$ on
  $G$}; its differential at $e$, again denoted $\mathrm{Ad}_g:\mathfrak g\to\mathfrak g$,
  is a Lie algebra isomorphism, the \emph{adjoint action of $G$ on $\mathfrak g$}.
\end{definition}

\begin{proposition}[$\mathrm{Ad}_g$ intertwines conjugation and the exponential map]
  \label{prop:pet-ch4-adjoint-exp-relation}
  \lean{PetersenLib.adjointExpRelation}
  \source{petersen:page-0156}
  \uses{def:pet-ch4-adjoint-action}
  \leanok
  $\mathrm{Ad}_g(\exp(tX))=\exp(t\,\mathrm{Ad}_g(X))$ for all $g\in G$, $X\in\mathfrak g$, $t\in\mathbb{R}$.
\end{proposition}
\begin{proof}
  \uses{prop:pet-ch4-adjoint-exp-relation}
  \leanok
  \sketch
  Conjugating the one-parameter subgroup $t\mapsto\exp(tX)$ by $g$ produces a
  one-parameter subgroup whose derivative at the identity is
  $\operatorname{Ad}_gX$. Uniqueness of the integral curve in
  \cref{prop:pet-ch4-matrix-exp-integral-curve} identifies it with
  $t\mapsto\exp(t\operatorname{Ad}_gX)$.
\end{proof}

\begin{proposition}[biinvariance criterion]
  \label{prop:pet-ch4-biinvariant-criterion}
  \lean{PetersenLib.biinvariantMetricCriterion}
  \source{petersen:page-0156,petersen:page-0157,petersen:page-0158}
  \uses{def:pet-ch4-adjoint-action, prop:pet-ch4-left-translation-isometry}
  \dcref{ch4:4.4.1}
  A left-invariant metric on a Lie group $G$ is biinvariant if and only if
  $\mathrm{Ad}_g:\mathfrak g\to\mathfrak g$ is a linear isometry for every $g\in G$.
\end{proposition}
\begin{proof}
  \uses{prop:pet-ch4-biinvariant-criterion}
  \sketch
  If the metric is biinvariant, conjugation $L_g\circ R_{g^{-1}}$ is an
  isometry, so its differential $\operatorname{Ad}_g$ is orthogonal. Conversely,
  $R_g=L_g\circ\operatorname{Ad}_{g^{-1}}$ at the identity; hence the hypothesis
  and \cref{prop:pet-ch4-left-translation-isometry} make $(DR_g)_e$ an
  isometry. For any $h$,
  $(DR_g)_h=(DR_{hg})_e\circ((DR_h)_e)^{-1}$, so every right translation is an
  isometry.
\end{proof}


\begin{proposition}[connection and curvature of a biinvariant metric]
  \label{prop:pet-ch4-biinvariant-curvature-formula}
  \lean{PetersenLib.biinvariantConnectionCurvature}
  \source{petersen:page-0158,petersen:page-0159}
  \uses{prop:pet-ch4-biinvariant-criterion, def:pet-ch4-adjoint-action}
  \dcref{ch4:4.4.2}
  On a Lie group $G$ with biinvariant metric $(\cdot,\cdot)$ and
  $X,Y,Z,W\in\mathfrak g$,
  \[
    \nabla_YX=\tfrac12[Y,X],\qquad R(X,Y)Z=-\tfrac14[[X,Y],Z],\qquad
    R(X,Y,Z,W)=\tfrac14\big([X,Y],[W,Z]\big).
  \]
  In particular, sectional curvature is always $\ge0$.
\end{proposition}
\begin{proof}
  \uses{prop:pet-ch4-biinvariant-curvature-formula}
  \sketch
  By \cref{prop:pet-ch4-biinvariant-criterion}, differentiating the orthogonal
  adjoint action shows that every $\operatorname{ad}_X$ is skew-adjoint. The
  Koszul formula therefore reduces to
  $2(\nabla_YX,Z)=([Y,X],Z)$, hence $\nabla_YX=\tfrac12[Y,X]$. Substitution in
  the curvature definition and the Jacobi identity give
  $R(X,Y)Z=-\tfrac14[[X,Y],Z]$. Skew-adjointness then yields
  $R(X,Y,Z,W)=\tfrac14([X,Y],[W,Z])$; for an orthonormal pair,
  $\operatorname{sec}(X,Y)=\tfrac14|[X,Y]|^2\ge0$.
\end{proof}


\begin{remark}[nonnegative curvature of biinvariant metrics]
  \label{rem:pet-ch4-biinvariant-nonneg-curvature}
  \lean{PetersenLib.biinvariantNonnegCurvatureRemark}
  \source{petersen:page-0159}
  \uses{prop:pet-ch4-biinvariant-curvature-formula}
  Lie groups with biinvariant Riemannian metrics always have nonnegative
  sectional curvature; with more work the curvature operator itself can be
  shown to be nonnegative.
\end{remark}

\subsection{Hyperbolic Space as a Lie Group}

\begin{example}[the upper half-plane group has constant curvature $-1$]
  \label{ex:pet-ch4-hyperbolic-lie-group}
  \lean{PetersenLib.hyperbolicPlaneAsLieGroup}
  \source{petersen:page-0159,petersen:page-0160}
  \uses{def:pet-ch4-left-invariant-metric-lie}
  On $G=\Big\{\left[\begin{smallmatrix}\alpha&\beta\\0&1\end{smallmatrix}\right]:\alpha>0,\beta\in\mathbb{R}\Big\}$
  (whose first row identifies $G$ with the upper half-plane), with Lie algebra
  $\mathfrak g=\Big\{\left[\begin{smallmatrix}a&b\\0&0\end{smallmatrix}\right]\Big\}$,
  $X=\left[\begin{smallmatrix}1&0\\0&0\end{smallmatrix}\right]$,
  $Y=\left[\begin{smallmatrix}0&1\\0&0\end{smallmatrix}\right]$, so $[X,Y]=Y$: declaring
  $X,Y$ orthonormal makes $G$ have constant curvature $-1$.
\end{example}
\begin{proof}
  \uses{ex:pet-ch4-hyperbolic-lie-group}
  \sketch
  With $[X,Y]=Y$, the Koszul formula gives
  $\nabla_XX=\nabla_XY=0$, $\nabla_YY=X$, and $\nabla_YX=-Y$. Therefore
  $R(X,Y)Y=-X$, so the unique sectional curvature in dimension two is $-1$.
\end{proof}

\begin{proposition}[the metric on the upper half-plane group is not biinvariant]
  \label{prop:pet-ch4-hyperbolic-lie-group-not-biinvariant}
  \lean{PetersenLib.hyperbolicLieGroupNotBiinvariant}
  \source{petersen:page-0160}
  \uses{ex:pet-ch4-hyperbolic-lie-group, prop:pet-ch4-biinvariant-criterion}
  For $g=\left[\begin{smallmatrix}\alpha&\beta\\0&1\end{smallmatrix}\right]\in G$,
  $\mathrm{Ad}_g\big(\left[\begin{smallmatrix}a&b\\0&0\end{smallmatrix}\right]\big)
  =\left[\begin{smallmatrix}a&-a\beta+b\alpha\\0&0\end{smallmatrix}\right]=aX+(-a\beta+b\alpha)Y$;
  this is not orthonormal-basis-preserving unless $\beta=0,\alpha=1$, so the
  metric on $G$ is not biinvariant (equivalently, the left-invariant fields
  $X,Y$ are not Killing fields).
\end{proposition}
\begin{proof}
  \uses{prop:pet-ch4-hyperbolic-lie-group-not-biinvariant}
  \sketch
  Direct conjugation sends the orthonormal basis to
  $\operatorname{Ad}_gX=X-\beta Y$ and
  $\operatorname{Ad}_gY=\alpha Y$. These vectors are orthonormal only when
  $\beta=0$ and $\alpha=1$. Thus the adjoint action is not generally isometric,
  and \cref{prop:pet-ch4-biinvariant-criterion} excludes biinvariance.
\end{proof}

\begin{remark}[higher-dimensional hyperbolic space as a Lie group, contrasted with spheres]
  \label{rem:pet-ch4-lie-group-spheres-contrast}
  \lean{PetersenLib.hyperbolicSpaceLieGroupHigherDim}
  \source{petersen:page-0161}
  \uses{prop:pet-ch4-hyperbolic-lie-group-not-biinvariant}
  \cref{ex:pet-ch4-hyperbolic-lie-group} generalizes: the upper half-space is
  in a natural way a Lie group with a left-invariant metric of constant
  curvature $-1$ in every dimension. This is in sharp contrast to the spheres,
  where only $S^3=SU(2)$ and $S^1=SO(2)$ are Lie groups.
\end{remark}

\subsection{Berger Spheres}

The $SU(2)$ frame calculations give the Berger curvature eigenvalues and their
degeneration under the Hopf fibration.

\begin{example}[general left-invariant metrics on $SU(2)$ diagonalize $\mathfrak{R}$]
  \label{ex:pet-ch4-general-berger-metric}
  \lean{PetersenLib.suTwoGeneralLeftInvariantMetric}
  \source{petersen:page-0161}
  \uses{prop:pet-ch4-biinvariant-criterion}
  On $SU(2)$, let $\lambda_1^{-1}X_1,\lambda_2^{-1}X_2,\lambda_3^{-1}X_3$ be
  orthonormal, where $[X_i,X_{i+1}]=2X_{i+2}$ (indices mod $3$). Then
  $\nabla_{X_i}X_i=0$ for each $i$, and
  \[
    \nabla_{X_i}X_{i+1}=\Big(\frac{\lambda_{i+2}^2+\lambda_{i+1}^2-\lambda_i^2}{\lambda_{i+2}^2}\Big)X_{i+2},
    \qquad
    \nabla_{X_{i+1}}X_i=\Big(\frac{-\lambda_{i+2}^2+\lambda_{i+1}^2-\lambda_i^2}{\lambda_{i+2}^2}\Big)X_{i+2},
  \]
  and $R(X_i,X_{i+1})X_{i+2}=0$ for all $i$, so all curvatures between three
  mutually distinct basis vectors vanish.
\end{example}
\begin{proof}
  \uses{ex:pet-ch4-general-berger-metric}
  \sketch
  Apply the Koszul formula in the orthogonal frame $X_1,X_2,X_3$, whose squared
  lengths are $\lambda_1^2,\lambda_2^2,\lambda_3^2$. It gives
  $\nabla_{X_i}X_i=0$ and the two displayed cyclic connection coefficients,
  paralleling \cref{prop:pet-ch4-biinvariant-curvature-formula} but retaining
  the unequal frame lengths. Substituting those coefficients into
  $R(X_i,X_{i+1})X_{i+2}$ shows that its three connection terms cancel, hence
  every curvature component involving three distinct frame directions is zero.
\end{proof}


\begin{example}[curvature of the Berger spheres]
  \label{ex:pet-ch4-berger-sphere-curvature}
  \lean{PetersenLib.bergerSphereCurvatureComputation}
  \source{petersen:page-0161,petersen:page-0162}
  \uses{ex:pet-ch4-general-berger-metric}
  Specializing \cref{ex:pet-ch4-general-berger-metric} to $\lambda_1=\varepsilon<1$,
  $\lambda_2=\lambda_3=1$ (the Berger sphere $g_\varepsilon$):
  \[
    \nabla_{X_1}X_2=(2-\varepsilon^2)X_3,\quad \nabla_{X_2}X_1=-\varepsilon^2X_3,\quad
    \nabla_{X_2}X_3=X_1,\quad \nabla_{X_3}X_2=-X_1,\quad
    \nabla_{X_3}X_1=\varepsilon^2X_2,\quad \nabla_{X_1}X_3=(\varepsilon^2-2)X_2,
  \]
  \[
    \mathfrak{R}(X_1\wedge X_2)=\varepsilon^2X_1\wedge X_2,\qquad
    \mathfrak{R}(X_3\wedge X_1)=\varepsilon^2X_3\wedge X_1,\qquad
    \mathfrak{R}(X_2\wedge X_3)=(4-3\varepsilon^2)X_2\wedge X_3,
  \]
  so all sectional curvatures lie in $[\varepsilon^2,4-3\varepsilon^2]$.
\end{example}
\begin{proof}
  \uses{ex:pet-ch4-berger-sphere-curvature}
  \sketch
  Substitute $\lambda_1=\varepsilon$ and $\lambda_2=\lambda_3=1$ into the
  cyclic connection formulas of \cref{ex:pet-ch4-general-berger-metric}. For
  example,
  \[
    R(X_1,X_2)X_2
    =-(2-\varepsilon^2)\nabla_{X_2}X_3-2\nabla_{X_3}X_2
    =\varepsilon^2X_1.
  \]
  The cyclic calculations give the other eigenvalues $\varepsilon^2$ and
  $4-3\varepsilon^2$. The three-distinct-index terms vanish by
  \cref{prop:pet-ch4-triple-vanishing-diagonalizes-operator}, so these are the
  curvature-operator eigenvalues; the range follows from
  \cref{prop:pet-ch4-curvature-operator-diagonal-bounds}.
\end{proof}


\begin{remark}[the Hopf-fibration limit $\varepsilon\to0$]
  \label{rem:pet-ch4-berger-limit}
  \lean{PetersenLib.bergerSphereLimitCurvature}
  \source{petersen:page-0162}
  \uses{ex:pet-ch4-berger-sphere-curvature}
  As $\varepsilon\to0$, $\mathrm{sec}(X_2,X_3)\to4$, the curvature of the base
  space $S^2(1/2)$ of the Hopf fibration.
\end{remark}

\begin{proposition}[adjoint action on the Berger-sphere frame]
  \label{prop:pet-ch4-berger-adjoint-action}
  \lean{PetersenLib.bergerSphereAdjointAction}
  \source{petersen:page-0162}
  \uses{ex:pet-ch4-general-berger-metric, prop:pet-ch4-biinvariant-criterion}
  For $g=\left[\begin{smallmatrix}z&w\\-\bar w&\bar z\end{smallmatrix}\right]\in SU(2)$,
  \[
    \mathrm{Ad}_gX_1=(|z|^2-|w|^2)X_1-2\mathrm{Re}(w\bar z)X_2-2\mathrm{Im}(w\bar z)X_3,
  \]
  \[
    \mathrm{Ad}_gX_2=2\,\mathrm{Im}(z\bar w)X_1+\mathrm{Re}(w^2+z^2)X_2+\mathrm{Im}(w^2+z^2)X_3,
  \]
  \[
    \mathrm{Ad}_gX_3=2\,\mathrm{Re}(z\bar w)X_1+\mathrm{Re}\big(i(z^2-w^2)\big)X_2+\mathrm{Im}\big(i(z^2-w^2)\big)X_3.
  \]
  If $X_1,X_2,X_3$ have equal length, $\mathrm{Ad}_g$ is an isometry for every
  $g$ (so the biinvariant metric on $SU(2)$ satisfies the criterion of
  \cref{prop:pet-ch4-biinvariant-criterion}); otherwise, generically, it is
  not.
\end{proposition}
\begin{proof}
  \uses{prop:pet-ch4-berger-adjoint-action}
  \sketch
  Conjugate each matrix $X_i$ by
  $g=\left[\begin{smallmatrix}z&w\\-\bar w&\bar z\end{smallmatrix}\right]$ and use
  $|z|^2+|w|^2=1$; expansion in the basis $X_1,X_2,X_3$ gives the three stated
  formulas. Their coefficient matrix is orthogonal when the three basis
  directions have equal length. If the lengths differ, its generic mixing of
  $X_1$ with $X_2,X_3$ fails to preserve the diagonal inner product.
\end{proof}

\section{Riemannian Submersions}

Horizontal and vertical decompositions control the connection and curvature of a
Riemannian submersion.

\subsection{Riemannian Submersions and Curvatures}

\begin{definition}[vertical and horizontal distributions, basic horizontal lift]
  \label{def:pet-ch4-vertical-horizontal-distribution}
  \lean{PetersenLib.verticalHorizontalDistribution}
  \source{petersen:page-0162,petersen:page-0163}
  For a Riemannian submersion $F:(\bar M,\bar g)\to(M,g)$, the \emph{vertical
  distribution} is $\mathcal V_{\bar p}=\ker DF_{\bar p}\subset T_{\bar p}\bar M$
  and the \emph{horizontal distribution} is $\mathcal H_{\bar p}=(\mathcal V_{\bar p})^\perp$;
  $DF:\mathcal H_{\bar p}\to T_pM$ is an isometry. For a vector field $X$ on
  $M$, its unique horizontal, $F$-related lift is the \emph{basic horizontal
  lift} $\bar X$. Every $v\in T_{\bar p}\bar M$ decomposes as $v=v^{\mathcal H}+v^{\mathcal V}$.
\end{definition}

\begin{proposition}[basic properties of vertical and basic horizontal fields]
  \label{prop:pet-ch4-submersion-basic-properties}
  \lean{PetersenLib.submersionBasicHorizontalProperties}
  \source{petersen:page-0163}
  \uses{def:pet-ch4-vertical-horizontal-distribution}
  \dcref{ch4:4.5.1}
  Let $V$ be a vertical vector field on $\bar M$ and $X,Y,Z$ vector fields on
  $M$ with basic horizontal lifts $\bar X,\bar Y,\bar Z$. Then:
  \begin{enumerate}
  \item $[V,\bar X]$ is vertical.
  \item $(L_V\bar g)(\bar X,\bar Y)=D_V\bar g(\bar X,\bar Y)=0$.
  \item $\bar g([\bar X,\bar Y],V)=2\bar g(\nabla_{\bar Y}\bar X,V)=-2\bar g(\nabla_{\bar X}\bar Y,V)=2\bar g(\nabla_VY,X)$.
  \item $\nabla_{\bar X}\bar Y=\overline{\nabla_XY}+\tfrac12[\bar X,\bar Y]^{\mathcal V}$.
  \end{enumerate}
\end{proposition}
\begin{proof}
  \uses{prop:pet-ch4-submersion-basic-properties}
  \sketch
  Since $V$ and $\bar X$ are related to $0$ and $X$, respectively,
  $DF[V,\bar X]=[0,X]=0$, proving (1). For (2), the bracket terms in
  $L_V\bar g(\bar X,\bar Y)$ vanish by (1), while
  $\bar g(\bar X,\bar Y)=g(X,Y)\circ F$ is constant along each fiber. Applying
  the Koszul formula to two basic horizontal fields and one vertical field
  gives all identities in (3). Its vertical projection is
  $\tfrac12[\bar X,\bar Y]^{\mathcal V}$; applying Koszul to three basic fields
  identifies the horizontal projection with $\overline{\nabla_XY}$, proving
  (4).
\end{proof}


\begin{definition}[the integrability tensor of a Riemannian submersion]
  \label{def:pet-ch4-integrability-tensor}
  \lean{PetersenLib.submersionIntegrabilityTensor}
  \source{petersen:page-0164}
  \uses{prop:pet-ch4-submersion-basic-properties}
  For basic horizontal $\bar X,\bar Y$, $[\bar X,\bar Y]^{\mathcal V}$ is
  tensorial and skew-symmetric in $\bar X,\bar Y$ (since
  $[\bar X,f\bar Y]^{\mathcal V}=f[\bar X,\bar Y]^{\mathcal V}+(D_{\bar X}f)\bar Y^{\mathcal V}=f[\bar X,\bar Y]^{\mathcal V}$,
  as $\bar Y$ is horizontal), defining the \emph{integrability tensor}
  $\mathcal H\times\mathcal H\to\mathcal V$; it measures the failure of
  Frobenius integrability of the horizontal distribution.
\end{definition}

\begin{example}[the Hopf map has non-integrable horizontal distribution]
  \label{ex:pet-ch4-hopf-integrability}
  \lean{PetersenLib.hopfMapIntegrabilityTensor}
  \source{petersen:page-0164}
  \uses{def:pet-ch4-integrability-tensor}
  \dcref{ch4:4.5.2}
  For the Hopf map $S^3(1)\to S^2(1/2)$, with orthonormal left-invariant frame
  $X_1,X_2,X_3$ satisfying $[X_i,X_{i+1}]=2X_{i+2}$: $X_1$ is vertical and
  $X_2,X_3$ are horizontal (though not basic); since $[X_2,X_3]=2X_1$, the
  horizontal distribution is not integrable.
\end{example}

\begin{theorem}[the O'Neill curvature formula]
  \label{thm:pet-ch4-oneill-formula}
  \lean{PetersenLib.oneillCurvatureFormula}
  \source{petersen:page-0164,petersen:page-0165}
  \uses{prop:pet-ch4-submersion-basic-properties, def:pet-ch4-integrability-tensor}
  \dcref{ch4:4.5.3}
  For a Riemannian submersion $F:(\bar M,\bar g)\to(M,g)$ with curvature
  tensors $R$ on $M$ and $\bar R$ on $\bar M$,
  \[
    g(R(X,Y)Y,X)=\bar g(\bar R(\bar X,\bar Y)\bar Y,\bar X)+\tfrac34\big|[\bar X,\bar Y]^{\mathcal V}\big|^2.
  \]
\end{theorem}
\begin{proof}
  \uses{thm:pet-ch4-oneill-formula}
  \sketch
  Choose commuting fields $X,Y,Z,H$ on the base. Their lifted brackets are
  vertical, and \cref{prop:pet-ch4-submersion-basic-properties}(4) splits each
  lifted covariant derivative into the lift of the base derivative plus half
  its vertical bracket. Expanding the curvature definition and applying part
  (3) gives
  \[
    \bar g(\bar R(\bar X,\bar Y)\bar Z,\bar H)=g(R(X,Y)Z,H)
    -\tfrac14\bar g([\bar Y,\bar Z],[\bar X,\bar H])
    +\tfrac14\bar g([\bar X,\bar Z],[\bar Y,\bar H])
    -\tfrac12\bar g([\bar X,\bar Y],[\bar H,\bar Z]).
  \]
  Taking $H=X$ and $Z=Y$ combines the bracket terms into
  $-\tfrac34|[\bar X,\bar Y]^{\mathcal V}|^2$ on the left-hand side, which is
  equivalent to the stated formula.
\end{proof}


\subsection{Riemannian Submersions and Lie Groups}

\begin{remark}[fiber-homogeneous submersions from isometric group actions]
  \label{rem:pet-ch4-fiber-homogeneous-submersions}
  \lean{PetersenLib.fiberHomogeneousSubmersionExamples}
  \source{petersen:page-0165}
  \uses{def:pet-ch4-vertical-horizontal-distribution}
  Many examples of nonnegative or positive curvature arise from a compact
  group $G$ acting freely by isometries on $(\bar M,\bar g)$ with
  $M=G\backslash\bar M$; such a submersion is \emph{fiber homogeneous}, $G$
  acting transitively on fibers. Examples: $\mathbb{CP}^n=S^{2n+1}/S^1$,
  $TS^n=(SO(n+1)\times\mathbb{R}^n)/SO(n)$, $SU(3)/T^2$.
\end{remark}

\begin{proposition}[an isometric group action with manifold quotient is a Riemannian submersion]
  \label{prop:pet-ch4-group-action-submersion}
  \lean{PetersenLib.groupActionInducesRiemannianSubmersion}
  \source{petersen:page-0165,petersen:page-0166}
  \uses{rem:pet-ch4-fiber-homogeneous-submersions}
  If $G$ acts by isometries on $(\bar M,\bar g)$ with $M=\bar M/G$ a manifold
  and $F:\bar M\to M$ the quotient map a submersion, then: for
  $\mathfrak x\in\mathfrak g$, the vector field
  $X|_{\bar p}=\tfrac{d}{dt}\big|_{t=0}\exp(t\mathfrak x)\cdot\bar p$ is a
  Killing field (the flow of $X$ is by isometries $F^t(\bar p)=\exp(t\mathfrak x)\cdot\bar p$),
  and $\mathcal V_{\bar p}=\mathrm{span}\{X|_{\bar p}:\mathfrak x\in\mathfrak g\}$;
  the action preserves $\mathcal V$, i.e.\ $Dx(\mathcal V_{\bar p})=\mathcal V_{x\bar p}$
  for $x\in G$, hence preserves the orthogonal complements $\mathcal H_{\bar p}$;
  defining the metric on $T_pM$ via the isomorphism $DF:\mathcal H_{\bar p}\to T_pM$
  is then well defined (independent of the choice of $\bar p\in F^{-1}(p)$),
  making $F$ a Riemannian submersion.
\end{proposition}
\begin{proof}
  \uses{prop:pet-ch4-group-action-submersion,prop:pet-ch4-adjoint-exp-relation}
  \sketch
  Each one-parameter subgroup acts isometrically, so its infinitesimal field is
  Killing and these fields span the vertical space. For $x\in G$,
  \cref{prop:pet-ch4-adjoint-exp-relation} gives
  \[
    Dx(X_{\mathfrak x}|_{\bar p})
    =X_{\operatorname{Ad}_x\mathfrak x}|_{x\bar p},
  \]
  so the action preserves the vertical distribution and, by orthogonality, the
  horizontal distribution. The isomorphisms
  $DF:\mathcal H_{\bar p}\to T_{F(\bar p)}M$ therefore induce the same inner
  product for every point in a fiber, defining the required quotient metric.
\end{proof}


\subsection{Complex Projective Space}

The Hopf quotient supplies the submersion metric on $\mathbb{CP}^n$ and its
sectional and Ricci curvature formulas.

\begin{definition}[the submersion metric on $\mathbb{CP}^n$]
  \label{def:pet-ch4-cpn-submersion-setup}
  \lean{PetersenLib.cpnSubmersionSetup}
  \source{petersen:page-0166,petersen:page-0167}
  \uses{prop:pet-ch4-group-action-submersion}
  $\mathbb{CP}^n=S^{2n+1}/S^1$ for the complex-scalar $S^1$-action on
  $S^{2n+1}\subset\mathbb{C}^{n+1}$. Writing
  $ds_{2n+1}^2=dr^2+\sin^2(r)ds_{2n-1}^2+\cos^2(r)d\theta^2$ and letting the
  $S^1$-action act separately on the $S^{2n-1}$ and $S^1$ factors,
  $\mathbb{CP}^n=[0,\pi/2]\times\big((S^{2n-1}\times S^1)/S^1\big)$, with metric
  $dr^2+\sin^2(r)(g+\cos^2(r)h)$ (in the splitting $ds_{2n-1}^2=h+g$ along and
  orthogonal to the residual Hopf circle). Let $V$ be the unit vertical field
  on $S^{2n+1}\subset\mathbb{C}^{n+1}$ tangent to the $S^1$ action.
\end{definition}

\begin{proposition}[the integrability tensor of the $\mathbb{CP}^n$ submersion]
  \label{prop:pet-ch4-cpn-integrability-tensor}
  \lean{PetersenLib.cpnIntegrabilityTensor}
  \source{petersen:page-0167}
  \uses{def:pet-ch4-cpn-submersion-setup, def:pet-ch4-integrability-tensor}
  For basic horizontal $\bar X,\bar Y$ on $S^{2n+1}$, $\tfrac12[\bar X,\bar Y]^{\mathcal V}=\bar g(\bar Y,i\bar X)V$,
  where $iV$ is the unit inward normal to $S^{2n+1}\subset\mathbb{C}^{n+1}$ and
  $\bar g$ the canonical Euclidean metric on $\mathbb{C}^{n+1}$; in particular the
  horizontal distribution (orthogonal to $V$) is invariant under
  multiplication by $i$, reflecting the complex structure of $\mathbb{CP}^n$.
\end{proposition}
\begin{proof}
  \uses{prop:pet-ch4-cpn-integrability-tensor}
  \sketch
  By \cref{prop:pet-ch4-submersion-basic-properties}(3),
  $\bar g(\tfrac12[\bar X,\bar Y],V)
  =\bar g(\nabla^{S^{2n+1}}_{\bar X}\bar Y,V)$. Compute this inner product in
  $\mathbb C^{n+1}$: metric compatibility and the parallel complex structure
  give
  \[
    \bar g(\nabla_{\bar X}\bar Y,V)
    =-\bar g(\bar Y,\nabla_{\bar X}V)
    =\bar g(\bar Y,i\bar X).
  \]
  Since $V$ spans the vertical distribution, this proves the formula. Orthogonality
  of $i$ and the fact that $iV$ is normal to the sphere show that $i$ preserves
  the horizontal space.
\end{proof}


\begin{theorem}[sectional curvature of $\mathbb{CP}^n$]
  \label{thm:pet-ch4-cpn-sectional-curvature}
  \lean{PetersenLib.cpnSectionalCurvature}
  \source{petersen:page-0167,petersen:page-0168}
  \uses{thm:pet-ch4-oneill-formula, prop:pet-ch4-cpn-integrability-tensor}
  For $X,Y$ orthonormal on $\mathbb{CP}^n$ with orthonormal horizontal lifts
  $\bar X,\bar Y$,
  \[
    \mathrm{sec}(X,Y)=1+3\big|\bar g(\bar Y,i\bar X)\big|^2\le4,
  \]
  with equality iff $\bar Y=\pm i\bar X$; in particular $\mathrm{sec}\ge1$
  everywhere on $\mathbb{CP}^n$.
\end{theorem}
\begin{proof}
  \uses{thm:pet-ch4-cpn-sectional-curvature}
  \sketch
  The total sphere has sectional curvature $1$ by
  \cref{prop:pet-ch4-euclidean-sphere-constant-curvature}. Combining
  \cref{thm:pet-ch4-oneill-formula} with
  \cref{prop:pet-ch4-cpn-integrability-tensor} gives
  \[
    \operatorname{sec}(X,Y)
    =1+\frac34\left|2\bar g(\bar Y,i\bar X)V\right|^2
    =1+3|\bar g(\bar Y,i\bar X)|^2.
  \]
  Cauchy--Schwarz bounds the final square by $1$, with equality exactly when
  $\bar Y=\pm i\bar X$.
\end{proof}

\begin{remark}[curvature-operator eigenbasis and eigenvalue range on $\mathbb{CP}^n$]
  \label{rem:pet-ch4-cpn-curvature-operator-eigenbasis}
  \lean{PetersenLib.cpnCurvatureOperatorEigenbasis}
  \source{petersen:page-0168}
  \uses{thm:pet-ch4-oneill-formula, thm:pet-ch4-cpn-sectional-curvature}
  Applying the O'Neill formula for the full curvature tensor to orthonormal
  vectors of the form $X,iX,Y,iY$ shows the curvature operator on $\mathbb{CP}^n$
  is not diagonalized by simple bivectors $E_i\wedge E_j$; instead it is
  diagonalized by
  $X\wedge iX\pm Y\wedge iY$, $X\wedge Y\pm iX\wedge iY$, $X\wedge iY\pm Y\wedge iX$,
  with eigenvalues in $[0,6]$.
\end{remark}

\begin{proposition}[$\mathbb{CP}^n$ is Einstein with constant $2n+2$]
  \label{prop:pet-ch4-cpn-einstein}
  \lean{PetersenLib.cpnEinsteinConstant}
  \source{petersen:page-0168}
  \uses{thm:pet-ch4-cpn-sectional-curvature}
  $\mathrm{Ric}(X,X)=2n+2$ for every unit $X$ on $\mathbb{CP}^n$.
\end{proposition}
\begin{proof}
  \uses{prop:pet-ch4-cpn-einstein}
  \sketch
  Extend $i\bar X$ to an orthonormal basis of $\bar X^\perp$. In the sectional
  curvature formula of \cref{thm:pet-ch4-cpn-sectional-curvature}, the complex
  line through $X$ contributes $4$, while each of the remaining $2n-2$
  orthogonal directions contributes $1$. Hence
  $\operatorname{Ric}(X,X)=4+(2n-2)=2n+2$.
\end{proof}

\subsection{Berger--Cheeger Perturbations}

\begin{definition}[the Berger--Cheeger construction $G\times_GM$]
  \label{def:pet-ch4-berger-cheeger-construction}
  \lean{PetersenLib.bergerCheegerConstruction}
  \source{petersen:page-0169}
  \uses{def:pet-ch4-left-invariant-metric-lie}
  Fix $(M,g)$ and a Lie group $G$ with right-invariant metric $(\cdot,\cdot)$
  acting by isometries on $M$. Then $G$ acts freely by isometries on
  $(G\times M,\,g_\lambda)$, $g_\lambda=\lambda(\cdot,\cdot)+g$ ($\lambda>0$), via
  $h\cdot(x,p)=(xh^{-1},hp)$. The quotient map $Q:G\times M\to M=(G\times M)/G$,
  $Q(x,p)=xp$, exhibits $M$ (diffeomorphic to the quotient) with a submersion
  metric $g_\lambda$, called the \emph{Berger--Cheeger perturbation} of $g$.
\end{definition}

\begin{proposition}[horizontal lift and length formula for the Berger--Cheeger metric]
  \label{prop:pet-ch4-berger-cheeger-horizontal-lift}
  \lean{PetersenLib.bergerCheegerHorizontalLift}
  \source{petersen:page-0169,petersen:page-0170}
  \uses{def:pet-ch4-berger-cheeger-construction}
  For $X\in T_eG$ generating the Killing field $X|_p=\tfrac d{dt}|_{t=0}\exp(tX)\cdot p$
  on $M$, the horizontal lift of $X|_p\in T_pM$ to $T_eG\times T_pM$ (for the
  action on $G\times M$) is
  \[
    \overline{X|_p}=\Big(\frac{|X|_p^2}{\lambda|X|^2+|X|_p^2}X,\ \frac{\lambda|X|^2}{\lambda|X|^2+|X|_p^2}X|_p\Big),
  \]
  with length $\big|\overline{X|_p}\big|_{g_\lambda}^2=\dfrac{\lambda|X|^2}{\lambda|X|^2+|X|_p^2}|X|_p^2\le|X|_p^2$,
  tending to $0$ as $\lambda\to0$ and to $|X|_p^2$ as $\lambda\to\infty$.
\end{proposition}
\begin{proof}
  \uses{prop:pet-ch4-berger-cheeger-horizontal-lift}
  \sketch
  At $(e,p)$ the infinitesimal vertical direction generated by $X$ is
  $(-X,X|_p)$. The vector
  $(|X|_p^2X,\lambda|X|^2X|_p)$ is $g_\lambda$-orthogonal to it, and
  its image under $DQ$ is
  $(\lambda|X|^2+|X|_p^2)X|_p$. Normalizing gives the displayed lift. Directly
  evaluating its product-metric norm gives
  \[
    |\overline{X|_p}|_{g_\lambda}^2
    =\frac{\lambda|X|^2|X|_p^2}{\lambda|X|^2+|X|_p^2},
  \]
  from which the inequality and both limits follow.
\end{proof}


\begin{remark}[curvature increase on the un-squeezed directions]
  \label{rem:pet-ch4-berger-cheeger-curvature-increase}
  \lean{PetersenLib.bergerCheegerCurvatureIncrease}
  \source{petersen:page-0170}
  \uses{prop:pet-ch4-berger-cheeger-horizontal-lift, thm:pet-ch4-oneill-formula}
  For $X,Y\in\mathcal H_p$ (already horizontal for the $G$-action on $M$, hence
  also horizontal for the action on $G\times M$),
  $\mathrm{sec}_{g_\lambda}(X,Y)\ge\mathrm{sec}_g(X,Y)$, the excess coming from
  the nonnegative integrability-tensor correction in
  \cref{thm:pet-ch4-oneill-formula} applied to the action on $G\times M$.
\end{remark}

\section{Exercises}

\begin{remark}[Schwarzschild curvature is not parallel]
  \label{rem:pet-ch4-ex-1}
  \lean{PetersenLib.exercise4_7_1}
  \source{petersen:page-0171}
  \uses{prop:pet-ch4-schwarzschild-ricci-flat}
  \dcref{ch4:4.7.1}
  Show that the Schwarzschild metric does not have parallel curvature tensor.
\end{remark}

\begin{remark}[Berger sphere curvature is not parallel]
  \label{rem:pet-ch4-ex-2}
  \lean{PetersenLib.exercise4_7_2}
  \source{petersen:page-0171}
  \uses{ex:pet-ch4-berger-sphere-curvature}
  \dcref{ch4:4.7.2}
  Show that the Berger spheres ($\varepsilon\ne1$) do not have parallel
  curvature tensor.
\end{remark}

\begin{remark}[isometries of projective spaces force $\nabla R=0$]
  \label{rem:pet-ch4-ex-3}
  \lean{PetersenLib.exercise4_7_3}
  \source{petersen:page-0171}
  \uses{def:pet-ch4-cpn-submersion-setup}
  \dcref{ch4:4.7.3}
  (1) Show $U(n+1)$ acts by isometries on $\mathbb{CP}^n$ (it acts by isometries on
  $S^{2n+1}(1)$ commuting with the quotient action). (2) Show that for each
  $p\in\mathbb{CP}^n$ there is an isometry $A_p\in\mathrm{Iso}_p$ with
  $DA_p|_p=-I$. (3) Use that isometries preserve $\nabla$ and $R$ to conclude
  $\nabla R=0$. (4) Repeat (1)--(3) for $\mathbb{HP}^n$ using the symplectic
  group $\mathrm{Sp}(n+1)$ of quaternionic matrices with $A^*A=I$.
\end{remark}

\begin{remark}[a more general Hessian criterion for a warped product]
  \label{rem:pet-ch4-ex-4}
  \lean{PetersenLib.exercise4_7_4}
  \source{petersen:page-0171,petersen:page-0172}
  \uses{thm:pet-ch4-brinkmann-warped-product}
  \dcref{ch4:4.7.4}
  If $(M,g)$ has a function $f$ with
  $\mathrm{Hess}\,f=\lambda(x)g+\mu(f)df^2$ for $\lambda:M\to\mathbb{R}$, $\mu:\mathbb{R}\to\mathbb{R}$,
  show that the metric is locally a warped product.
\end{remark}

\begin{remark}[the conformal factor is $\Delta f/\dim M$]
  \label{rem:pet-ch4-ex-5}
  \lean{PetersenLib.exercise4_7_5}
  \source{petersen:page-0172}
  \uses{thm:pet-ch4-brinkmann-warped-product}
  \dcref{ch4:4.7.5}
  Show that if $\mathrm{Hess}\,f=\lambda g$ then $\lambda=\Delta f/\dim M$.
\end{remark}

\begin{remark}[warped product from an eigenvector Hessian with $\le2$ eigenvalues]
  \label{rem:pet-ch4-ex-6}
  \lean{PetersenLib.exercise4_7_6}
  \source{petersen:page-0172}
  \uses{thm:pet-ch4-brinkmann-warped-product}
  \dcref{ch4:4.7.6}
  For $f$ on $(M,g)$ with $\nabla f\ne0$ an eigenvector of $S(X)=\nabla_X\nabla f$:
  show that if $S$ has $\le2$ eigenvalues the metric is locally a warped
  product.
\end{remark}

\begin{remark}[the $A$-tensors of a Riemannian submersion]
  \label{rem:pet-ch4-ex-7}
  \lean{PetersenLib.exercise4_7_7}
  \source{petersen:page-0172}
  \uses{def:pet-ch4-integrability-tensor, prop:pet-ch4-submersion-basic-properties}
  \dcref{ch4:4.7.7}
  Define $A_{\bar X}\bar Y=(\bar\nabla_{\bar X}\bar Y)^{\mathcal V}$,
  $A_{\bar X}V=(\bar\nabla_{\bar X}V)^{\mathcal H}$, extended by $A_V=0$, and let
  $T_V$ be the second-fundamental-form tensor of the fibers (extended by
  $T_{\bar X}=0$), giving $(1,2)$-tensors on $\bar M$. (1) Show both
  $A$-tensors are tensorial. (2) Show
  $A_{\bar X}\bar Y=\tfrac12[\bar X,\bar Y]^{\mathcal V}$. (3) Show
  $\bar g(A_{\bar X}\bar Y,V)=-\bar g(\bar Y,A_{\bar X}V)$. (4) Show
  $(\nabla_VA)_W=-A_{T_{V,W}}$ and $(\nabla_{\bar X}A)_W=-A_{A_{\bar X}W}$.
  (5) Show $(\nabla_{\bar X}T)_{\bar Y}=-T_{A_{\bar X}\bar Y}$ and
  $(\nabla_VT)_{\bar Y}=-T_{T_{V,\bar Y}}$. (6) Show
  $\bar g\big((\nabla_UA)_{\bar X}V,W\big)=\bar g(T_UV,A_{\bar X}W)-\bar g(T_UW,A_{\bar X}V)$.
\end{remark}

\begin{remark}[horizontal and mixed curvature formulas]
  \label{rem:pet-ch4-ex-8}
  \lean{PetersenLib.exercise4_7_8}
  \source{petersen:page-0172}
  \uses{rem:pet-ch4-ex-7, thm:pet-ch4-oneill-formula}
  \dcref{ch4:4.7.8}
  Using the $A,T$-tensors of \cref{rem:pet-ch4-ex-7}, show
  $\bar R(\bar Y,\bar X,\bar X,\bar Y)=R(Y,X,X,Y)-3|A_{\bar X}\bar Y|^2$ and
  $\bar R(V,\bar X,\bar X,V)=\bar g\big((\nabla_{\bar X}T)_V,\bar X\big)+|A_{\bar X}V|^2-|T_V\bar X|^2$;
  compare the second formula to the radial curvature equation.
\end{remark}

\begin{remark}[curvature, Einstein condition, and Weyl tensor of a product]
  \label{rem:pet-ch4-ex-9}
  \lean{PetersenLib.exercise4_7_9}
  \source{petersen:page-0172,petersen:page-0173}
  \dcref{ch4:4.7.9}
  Let $(M,g)=(M_1\times M_2,g_1+g_2)$. (1) Show $R=R_1+R_2$ (pulled back). (2)
  Assume $(M_i,g_i)$ has constant curvature $c_i$; show $R=c_1g_1\circ g_1+c_2g_2\circ g_2$.
  (3) Show $(M,g)$ is Einstein iff $(n_1-1)c_1=(n_2-1)c_2$, $n_i=\dim M_i$. (4)
  Show the Weyl tensor vanishes if $c_1=-c_2$, $n_1=1$, or $n_2=1$. (5) Show
  the Weyl tensor does not vanish if none of these hold.
\end{remark}

\begin{remark}[constant curvature and the double warped-product representation of hyperbolic space]
  \label{rem:pet-ch4-ex-10}
  \lean{PetersenLib.exercise4_7_10}
  \source{petersen:page-0173}
  \uses{thm:pet-ch4-rotsym-curvature-operator}
  \dcref{ch4:4.7.10}
  Let $(M^n,g)=(I\times N,dr^2+\rho^2(r)g_N)$ have constant curvature $k$. (1)
  Show $(N^{n-1},\rho^2(r)g_N)$ has constant curvature $k+(\dot\rho/\rho)^2$ if
  $n>2$. (2) Show explicitly that hyperbolic space is a warped product over
  both hyperbolic space and Euclidean space.
\end{remark}

\begin{remark}[Einstein metric via a rescaling function]
  \label{rem:pet-ch4-ex-11}
  \lean{PetersenLib.exercise4_7_11}
  \source{petersen:page-0173}
  \uses{prop:pet-ch4-rotsym-ricci-scalar}
  \dcref{ch4:4.7.11}
  Let $(N^{n-1},g_N)$ have $\mathrm{Ric}=\tfrac{n-2}2\lambda g_N$, $\lambda<0$.
  Find $\rho:\mathbb{R}\to(0,\infty)$ so that $(M^n,g)=(\mathbb{R}\times N,dr^2+\rho^2(r)g_N)$
  is Einstein with $\mathrm{Ric}=\lambda g$.
\end{remark}

\begin{remark}[curvature, Weyl tensor, and Schouten tensor of a constant-curvature-base warped product]
  \label{rem:pet-ch4-ex-12}
  \lean{PetersenLib.exercise4_7_12}
  \source{petersen:page-0173}
  \uses{thm:pet-ch4-rotsym-curvature-operator}
  \dcref{ch4:4.7.12}
  Let $(N^{n-1},g_N)$ have constant curvature $c$, $n>2$, and
  $(M,g)=(I\times N,dr^2+\rho^2(r)g_N)$. (1) Show
  $R=\tfrac{c-\dot\rho^2}{\rho^2}g_r\circ g_r-2\tfrac{\ddot\rho}\rho dr^2\circ g_r
  =\tfrac{c-\dot\rho^2}{\rho^2}g\circ g-2\big(\tfrac{\ddot\rho}\rho+\tfrac{c-\dot\rho^2}{\rho^2}\big)dr^2\circ g$.
  (2) Show the Weyl tensor vanishes. (3) Show directly that the Schouten
  tensor $P$ satisfies $(\nabla_XP)(Y,Z)=(\nabla_YP)(X,Z)$.
\end{remark}

\begin{remark}[stereographic projection and the conformal models]
  \label{rem:pet-ch4-ex-13}
  \lean{PetersenLib.exercise4_7_13}
  \source{petersen:page-0173,petersen:page-0174}
  \uses{prop:pet-ch4-sphere-stereographic-conformal, prop:pet-ch4-poincare-disc-model}
  \dcref{ch4:4.7.13}
  For $S(x)=-e_{n+1}+\lambda(x)(e_{n+1}+(x,0))$, $x\in\mathbb{R}^n$, stereographic
  projection from $-e_{n+1}$ onto a hypersurface $M$: (1) for $M=S^n(1)$, show
  $\lambda(1+|x|^2)=2$ and $S$ is conformal with
  $g_{S^n(1)}=\tfrac4{(1+|x|^2)^2}g_{\mathbb{R}^n}$; (2) for $M=H^n(1)\subset\mathbb{R}^{n,1}$,
  show $\lambda(1-|x|^2)=2$ and $S$ is conformal with the Poincaré disc metric
  $\tfrac4{(1-|x|^2)^2}g_{\mathbb{R}^n}$.
\end{remark}

\begin{remark}[conformal change formulas and Weyl invariance]
  \label{rem:pet-ch4-ex-14}
  \lean{PetersenLib.exercise4_7_14}
  \source{petersen:page-0174}
  \dcref{ch4:4.7.14}
  Let $\tilde g=e^{2\psi}g$. (1) Show
  $\tilde\nabla_XY=\nabla_XY+(D_X\psi)Y+(D_Y\psi)X-g(X,Y)\nabla\psi$. (2) Show
  $e^{-2\psi}\tilde R=R-\big(2\,\mathrm{Hess}\,\psi-2(d\psi)^2+|d\psi|^2g\big)\circ g$.
  (3) For orthonormal $X,Y$, show
  $e^{2\psi}\widetilde{\mathrm{sec}}(X,Y)=\mathrm{sec}(X,Y)-\mathrm{Hess}\,\psi(X,X)-\mathrm{Hess}\,\psi(Y,Y)+(D_X\psi)^2+(D_Y\psi)^2-|d\psi|^2$.
  (4) Show
  $\widetilde{\mathrm{Ric}}=\mathrm{Ric}-(n-2)(\mathrm{Hess}\,\psi-d\psi^2)-(\Delta\psi+(n-2)|d\psi|^2)g$.
  (5) Show $e^{2\psi}\widetilde{\mathrm{scal}}=\mathrm{scal}-2(n-1)\Delta\psi-(n-1)(n-2)|d\psi|^2$.
  (6) Using exercise 3.4.25, show $e^{-2\psi}\tilde W=W$ (conformal invariance
  of the Weyl tensor, Weyl).
\end{remark}

\begin{remark}[explicit algebraic rewriting of the scalar-flat and Schwarzschild metrics]
  \label{rem:pet-ch4-ex-15}
  \lean{PetersenLib.exercise4_7_15}
  \source{petersen:page-0175}
  \uses{ex:pet-ch4-scalar-flat-tautological-bundle, prop:pet-ch4-schwarzschild-ricci-flat}
  \dcref{ch4:4.7.15}
  Show
  $\big(\tfrac14\rho_0^2+r^{2-n}\big)^{\frac{n-2}2}g_{\mathbb{R}^n}
  =\tfrac1{1-(\rho_0/\rho)^{n-2}}d\rho^2+\rho^2ds_{n-1}^2$, the scalar-flat
  metric of \cref{ex:pet-ch4-scalar-flat-tautological-bundle}; use this to
  rewrite the Schwarzschild metric of \cref{prop:pet-ch4-schwarzschild-ricci-flat}
  algebraically in terms of the flat metric on $\mathbb{R}^{n-1}$ and a conformal
  factor.
\end{remark}

\begin{remark}[the static Einstein equations]
  \label{rem:pet-ch4-ex-16}
  \lean{PetersenLib.exercise4_7_16}
  \source{petersen:page-0175,petersen:page-0176}
  \dcref{ch4:4.7.16}
  For $(M,g)=(N\times\mathbb{R},g_N+w^2dt^2)$, $w:N\to(0,\infty)$, $\dim N=n-1$: (1)
  show $\nabla^N_XY=\nabla^M_XY$, $R^N(X,Y)Z=R^M(X,Y)Z$ for $X,Y,Z$ tangent to
  $N$, and conclude $\mathrm{Ric}^M(X,\partial_t)=0$. (2) Show $|\partial_t|^2=w^2$.
  (3) Show $\nabla^M_{\partial_t}\partial_t=-w\nabla w$ and
  $\nabla^M_X\partial_t=\nabla^M_{\partial_t}X=\tfrac1w(D_Xw)\partial_t$. (4)
  Show $R^M(X,\partial_t)\partial_t=-w\nabla_X\nabla w$,
  $\mathrm{Ric}^M(\partial_t,\partial_t)=-w\Delta w$,
  $\mathrm{Ric}^M(X,X)=\mathrm{Ric}^N(X,X)-\tfrac1w\mathrm{Hess}\,w(X,X)$. (5) Show
  $\mathrm{Ric}^M=\lambda g$ iff $\mathrm{Ric}^N-\tfrac1w\mathrm{Hess}\,w=\lambda g_N$
  and $w\Delta w+\lambda w^2=0$, iff
  $\mathrm{Ric}^N-\tfrac1w\mathrm{Hess}\,w=\lambda g_N$ and $\mathrm{scal}^N=(n-2)\lambda$.
\end{remark}

\begin{remark}[local conformal flatness and isothermal coordinates]
  \label{rem:pet-ch4-ex-17}
  \lean{PetersenLib.exercise4_7_17}
  \source{petersen:page-0176}
  \uses{prop:pet-ch4-snk-constant-curvature}
  \dcref{ch4:4.7.17}
  A metric is \emph{locally conformally flat} if near every point
  $g=e^{-2\psi}\big((dx^1)^2+\dots+(dx^n)^2\big)$. (1) Show the space forms
  $S^n_k$, $dr^2+\mathrm{sn}_k^2(r)ds_{n-1}^2$, are locally conformally flat. (2)
  Show a locally conformally flat Einstein metric has constant curvature. (3)
  For $\dim M=2$ (Gauss's isothermal coordinates): (a) if $du\ne0$ on
  $O\subset M$, show there is, up to sign, a unique $1$-form
  $\omega=i_{\nabla u}\mathrm{vol}_g$ with $|du|=|\omega|$ and
  $d(du,\omega)=0$; (b) show $d\omega=(\Delta_gu)\mathrm{vol}_g$; (c) show
  isothermal coordinates exist near $p$ whenever there is a harmonic $u$ near
  $p$ with $du|_p\ne0$.
\end{remark}

\begin{remark}[local conformal flatness via the Schouten tensor]
  \label{rem:pet-ch4-ex-18}
  \lean{PetersenLib.exercise4_7_18}
  \source{petersen:page-0176,petersen:page-0177}
  \dcref{ch4:4.7.18}
  Let $\dim M=n>2$. (1) Show $g$ is locally conformally flat iff $W=0$ and
  locally there is $\psi$ with
  $P=2\,\mathrm{Hess}\,\psi-2(d\psi)^2+|d\psi|^2g$ (the condition $W=0$ being
  redundant for $n=3$). (2) Show that if $g$ is locally conformally flat then
  $(\nabla_XP)(Y,Z)=(\nabla_YP)(X,Z)$.
\end{remark}

\begin{remark}[the converse integrability argument]
  \label{rem:pet-ch4-ex-19}
  \lean{PetersenLib.exercise4_7_19}
  \source{petersen:page-0176,petersen:page-0177}
  \uses{rem:pet-ch4-ex-18}
  \dcref{ch4:4.7.19}
  Assume $\dim M=n>2$, $W=0$, and $(\nabla_XP)(Y,Z)=(\nabla_YP)(X,Z)$. (1) Show
  that if $\nabla\omega=\tfrac12P+\omega^2-\tfrac12|\omega|^2g$ for a $1$-form
  $\omega$, then locally $\omega=d\psi$ and
  $P=2\,\mathrm{Hess}\,\psi-2(d\psi)^2+|\nabla\psi|^2g$. (2) Using the Ricci
  identity $\nabla^2_{X,Y}\omega-\nabla^2_{Y,X}\omega=R_{X,Y}\omega$, show that
  the same hypothesis on $\nabla\omega$ gives
  $(\nabla^2_{X,Y}\omega)(Z)=\tfrac12(\nabla_XP)(Y,Z)+(\nabla_X\omega)(Y)\omega(Z)+\omega(Y)(\nabla_X\omega)(Z)-g(\nabla_X\omega,\omega)g(Y,Z)$.
  (3) Deduce $\nabla^2_{X,Y}\omega-\nabla^2_{Y,X}\omega=(P\circ g)(X,Y,V,Z)$ for
  $V$ dual to $\omega$. (4) Using $R=P\circ g$, show
  $(R_{X,Y}\omega)(Z)=(P\circ g)(X,Y,V,Z)$. (5) Conclude the integrability
  conditions for $\omega$ hold, so the manifold is locally conformally flat.
\end{remark}

\begin{remark}[conformal flatness of $N^2\times\mathbb{R}$]
  \label{rem:pet-ch4-ex-20}
  \lean{PetersenLib.exercise4_7_20}
  \source{petersen:page-0177}
  \dcref{ch4:4.7.20}
  For $(N^2\times\mathbb{R},g_N+g_{\mathbb{R}})$: (1) show
  $P_{N\times\mathbb{R}}=\tfrac{\mathrm{scal}_N}2(g_N-g_{\mathbb{R}})$; (2) show this product
  metric is conformally flat iff $\mathrm{scal}_N$ is constant.
\end{remark}

\begin{remark}[explicit conformal factor for constant curvature]
  \label{rem:pet-ch4-ex-21}
  \lean{PetersenLib.exercise4_7_21}
  \source{petersen:page-0178}
  \uses{rem:pet-ch4-ex-19}
  \dcref{ch4:4.7.21}
  Let $(M^n,g)$, $n>2$, have constant curvature $k$. (1) Use exercise 4.7.19
  to show the metric is locally conformally flat. (2) If
  $g=e^{-2\psi}\big((dx^1)^2+\dots+(dx^n)^2\big)$, show
  $2e^\psi\partial_i\partial_j\psi=\big(k+\sum_k(\partial_ke^\psi)^2\big)\delta_{ij}$.
  (3) Show $e^\psi=a+\sum_ib_ix^i+c\sum_i(x^i)^2$ with $k=4ac-\sum_ib_i^2$.
\end{remark}

\begin{remark}[the Heisenberg group]
  \label{rem:pet-ch4-ex-22}
  \lean{PetersenLib.exercise4_7_22}
  \source{petersen:page-0178}
  \uses{def:pet-ch4-left-invariant-metric-lie}
  \dcref{ch4:4.7.22}
  For the Heisenberg group $G$ of upper unitriangular $3\times3$ matrices with
  Lie algebra basis $X,Y,Z$: (1) show the only nonzero bracket is $[X,Y]=Z$;
  declare $X,Y,Z$ orthonormal for a left-invariant metric. (2) Show
  $\mathrm{Ric}$ has both positive and negative eigenvalues. (3) Show
  $\mathrm{scal}$ is constant. (4) Show $\mathrm{Ric}$ is not parallel.
\end{remark}

\begin{remark}[the Eguchi--Hanson metric is Ricci flat]
  \label{rem:pet-ch4-ex-23}
  \lean{PetersenLib.exercise4_7_23}
  \source{petersen:page-0178,petersen:page-0179}
  \uses{ex:pet-ch4-general-berger-metric}
  \dcref{ch4:4.7.23}
  For $dr^2+\rho^2(r)\big(\phi^2(r)(\sigma^1)^2+(\sigma^2)^2+(\sigma^3)^2\big)$
  (in the coframe dual to the left-invariant frame $X_1,X_2,X_3$ of
  \cref{ex:pet-ch4-general-berger-metric} on $S^3$):
  (1) show that $\dot\rho=\phi$, $\dot\rho^2=1-k\rho^{-4}$,
  $\rho(0)=k^{1/4}$, $\dot\rho(0)=0$, $\phi(0)=0$, $\dot\phi(0)=2$ give a family
  of Ricci-flat metrics on $TS^2$. (2) Show $\rho(r)\sim r$, $\dot\rho(r)\sim1$,
  and (a further derivative) $\sim2kr^{-5}$ as $r\to\infty$; conclude
  curvatures are $O(r^{-6})$ and the metric looks like
  $(0,\infty)\times\mathbb{RP}^3=(0,\infty)\times SO(3)$ at infinity, that
  rescaling corresponds to changing $k$, and that this is (up to scale) the
  unique Ricci-flat \emph{Eguchi--Hanson metric}.
\end{remark}

\begin{remark}[the Kähler/Hermitian structure and scalar-flat bundle metrics]
  \label{rem:pet-ch4-ex-24}
  \lean{PetersenLib.exercise4_7_24}
  \source{petersen:page-0179}
  \uses{rem:pet-ch4-ex-23}
  \dcref{ch4:4.7.24}
  For $dr^2+\rho^2(r)\big(\phi^2(r)(\sigma^1)^2+(\sigma^2)^2+(\sigma^3)^2\big)$
  and the $(1,1)$-tensor represented (in the orthonormal frame) by the block
  matrix $\left[\begin{smallmatrix}0&-1&0&0\\1&0&0&0\\0&0&0&-1\\0&0&1&0\end{smallmatrix}\right]$,
  which defines a Hermitian structure: (1) show this structure is Kähler
  (parallel) iff $\dot\rho=\phi$. (2) Find the scalar curvature of such
  metrics. (3) Show there are scalar-flat metrics on all $2$-dimensional
  vector bundles over $S^2$: the one on $TS^2$ is the Eguchi--Hanson metric,
  and the one on $S^2\times\mathbb{R}^2$ is the Schwarzschild metric.
\end{remark}

\begin{remark}[nonpositive Ricci curvature on the tautological bundle]
  \label{rem:pet-ch4-ex-25}
  \lean{PetersenLib.exercise4_7_25}
  \source{petersen:page-0179}
  \uses{ex:pet-ch4-scalar-flat-tautological-bundle}
  \dcref{ch4:4.7.25}
  Show that $\tau(\mathbb{RP}^{n-1})$ admits rotationally symmetric metrics
  $dr^2+\rho^2(r)ds_{n-1}^2$ with $\rho(r)=r$ for $r>1$ and nonpositive Ricci
  curvature, i.e.\ the Euclidean metric can be topologically perturbed to have
  nonpositive Ricci curvature (though not, similarly, to have nonnegative
  scalar curvature or nonpositive sectional curvature -- examine rotationally
  symmetric metrics on $\mathbb{R}^n$ and $\tau(\mathbb{RP}^{n-1})$).
\end{remark}

\begin{remark}[orthogonal coordinates]
  \label{rem:pet-ch4-ex-26}
  \lean{PetersenLib.exercise4_7_26}
  \source{petersen:page-0179,petersen:page-0180}
  \dcref{ch4:4.7.26}
  A manifold admits \emph{orthogonal coordinates} at $p$ if in some coordinate
  neighborhood $g_{ij}=0$ for $i\ne j$. Show such coordinates always exist in
  dimension $2$ (and, though not proved here, dimension $3$), but may fail to
  exist in dimension $>3$; for a counterexample, show that in orthogonal
  coordinates the curvatures $R^i_{jk\ell}=0$ whenever all indices are
  distinct.
\end{remark}

\begin{remark}[nonvanishing of the Schwarzschild and Eguchi--Hanson Weyl tensors]
  \label{rem:pet-ch4-ex-27}
  \lean{PetersenLib.exercise4_7_27}
  \source{petersen:page-0180}
  \uses{prop:pet-ch4-schwarzschild-ricci-flat, rem:pet-ch4-ex-23}
  \dcref{ch4:4.7.27}
  Show the Weyl tensors for the Schwarzschild metric and the Eguchi--Hanson
  metric are nonzero.
\end{remark}

\begin{remark}[the Hodge star and the block decomposition of the curvature operator]
  \label{rem:pet-ch4-ex-28}
  \lean{PetersenLib.exercise4_7_28}
  \source{petersen:page-0180,petersen:page-0181}
  \uses{thm:pet-ch4-doubly-warped-curvature, ex:pet-ch4-general-berger-metric, prop:pet-ch4-cpn-einstein}
  \dcref{ch4:4.7.28}
  Let $(M^4,g)$ be oriented, and $*:\Lambda^2TM\to\Lambda^2TM$ the Hodge star,
  $*(e_1\wedge e_2)=e_3\wedge e_4$ for oriented orthonormal $e_1,\dots,e_4$. (1)
  Show $*$ is a well-defined endomorphism with $**=I$ (and in dimension $2$,
  $*:TM\to TM$ satisfies $**=-I$, an almost complex structure). (2) Show
  $e_1\wedge e_2\pm e_3\wedge e_4$, $e_1\wedge e_3\pm e_4\wedge e_2$,
  $e_1\wedge e_4\pm e_2\wedge e_3$ lie in the $\pm1$ eigenspaces
  $\Lambda^\pm TM$ of $*$. (3) Deduce the block decomposition
  $\mathfrak R=\left[\begin{smallmatrix}A&D\\B&C\end{smallmatrix}\right]$ with
  $A=W^++\tfrac{\mathrm{scal}}{12}I$, $C=W^-+\tfrac{\mathrm{scal}}{12}I$,
  $D=B^*$; find these decompositions for the doubly warped metrics
  $dr^2+\rho^2(r)d\theta^2+\phi^2(r)ds_2^2$ on $I\times S^1\times S^2$ and
  $dr^2+\rho^2(r)\big(\phi^2(r)(\sigma^1)^2+(\sigma^2)^2+(\sigma^3)^2\big)$ on
  $I\times S^3$. (4) Find the curvature operators for the Schwarzschild
  metric, the Eguchi--Hanson metric, $S^2\times S^2$, $S^4$, and $\mathbb{CP}^2$.
  (5) Show $(M,g)$ is Einstein iff $B=0$ iff $\mathrm{sec}(\pi)=\mathrm{sec}(\pi^\perp)$
  for every plane $\pi$ and its orthogonal complement $\pi^\perp$.
\end{remark}
